\documentclass[11pt,reqno]{amsart}
\usepackage[margin=1.15in]{geometry}
\usepackage{amsmath,amssymb,amsthm,mathtools}
\usepackage{enumitem}
\usepackage{booktabs}
\usepackage{xcolor}
\usepackage[colorlinks=true,linkcolor=blue!50!black,citecolor=green!40!black,urlcolor=blue!50!black]{hyperref}

% ---------------------------------------------------------------------------
% Theorem environments
% ---------------------------------------------------------------------------
\newtheorem{theorem}{Theorem}[section]
\newtheorem{proposition}[theorem]{Proposition}
\newtheorem{lemma}[theorem]{Lemma}
\newtheorem{corollary}[theorem]{Corollary}
\newtheorem{fact}[theorem]{Fact}
\newtheorem{mainthm}{Theorem}
\renewcommand{\themainthm}{\Alph{mainthm}}
\theoremstyle{definition}
\newtheorem{definition}[theorem]{Definition}
\newtheorem{example}[theorem]{Example}
\theoremstyle{remark}
\newtheorem{remark}[theorem]{Remark}

% ---------------------------------------------------------------------------
% Macros
% ---------------------------------------------------------------------------
\newcommand{\R}{\mathbb{R}}
\newcommand{\N}{\mathbb{N}}
\newcommand{\Pp}{\mathcal{P}}
\newcommand{\Dom}{\mathcal{D}}
\newcommand{\Par}{\mathcal{O}}
\newcommand{\Lcl}{\mathcal{L}}
\newcommand{\Jet}{\mathrm{J}}
\newcommand{\Diff}{\operatorname{Diff}}
\newcommand{\Hom}{\operatorname{Hom}}
\newcommand{\End}{\operatorname{End}}
\newcommand{\tr}{\operatorname{tr}}
\newcommand{\id}{\mathrm{id}}
\newcommand{\dmu}{\,d\mu}
\newcommand{\ip}[2]{\langle #1,#2\rangle}
\newcommand{\norm}[1]{\lVert #1\rVert}
\newcommand{\abs}[1]{\lvert #1\rvert}
\newcommand{\Lap}{\nabla^{*}\nabla}
\newcommand{\hatP}{\widehat{P}}
\newcommand{\todo}[1]{\textcolor{red!70!black}{[#1]}}

% ---------------------------------------------------------------------------
% Blueprint annotations (leanblueprint-compatible). They are invisible in the
% PDF; scripts/blueprint.py reads them to build blueprint/map.md.
%   \lean{Decl1, Decl2}  Lean declaration(s) formalizing this item
%   \uses{label1, ...}   items this statement / proof depends on
%   \source{text}        where an imported Fact comes from
% ---------------------------------------------------------------------------
\newcommand{\lean}[1]{}
\newcommand{\uses}[1]{}
\newcommand{\leanok}{}
\newcommand{\source}[1]{}

\numberwithin{equation}{section}

\title[Short-time existence for higher-order parabolic systems]{Short-time existence, uniqueness and smooth dependence for nonlinear parabolic equations of arbitrary even order on vector bundles}
\author{[Authors]}
\date{Working draft, \today}

\begin{document}

\begin{abstract}
We prove that a fully nonlinear evolution equation
$\partial_t u = F(t,u,\nabla u,\dots,\nabla^{2p}u)$
for sections of a vector bundle over a closed Riemannian manifold has, for every smooth
initial datum at which the equation is parabolic, a unique smooth solution for a short time, and
that this solution depends smoothly on the initial datum. Parabolicity is expressed by a pointwise
algebraic \emph{Legendre condition} on the linearization. It covers the quasilinear equations of
product type treated by Mantegazza and Martinazzi, polyharmonic flows of every order, the
Ricci--DeTurck flow and parabolic Monge--Amp\`ere equations.

The proof is self-contained modulo standard facts about Sobolev spaces on closed manifolds, and it
is organized with a machine formalization in mind:
\begin{itemize}[nosep]
\item all function spaces have integer order;
\item no trace, interpolation or fractional spaces occur;
\item linear theory is needed only for operators with smooth coefficients;
\item regularity of solutions follows from the implicit function theorem via Angenent's parameter
  trick, not from tame estimates.
\end{itemize}
\end{abstract}

\maketitle

\noindent\emph{Status.} This is a working draft intended to serve as the mathematical blueprint of a
Lean formalization. Results labelled ``Fact'' are standard and are imported rather than proved here;
Section~\ref{sec:formal} lists them together with the dependency structure of the argument.
Places where the draft still needs checking are marked \todo{like this}.

\tableofcontents

% ===========================================================================
\section{Introduction}\label{sec:intro}
% ===========================================================================

Let $(M,g)$ be a closed Riemannian manifold and $V\to M$ a vector bundle equipped with a fibre metric
and a compatible connection. We study evolution equations
\begin{equation}\label{eq:intro}
  \partial_t u = F\bigl(t,u,\nabla u,\dots,\nabla^{2p}u\bigr),\qquad u(\cdot,0)=u_0,
\end{equation}
for time-dependent sections $u$ of $V$, where $p\ge 1$ and $F$ is smooth. Equations of this form
occur throughout geometric analysis:
\begin{itemize}[nosep]
\item order two ($p=1$): the harmonic map heat flow and the Ricci--DeTurck flow;
\item order four: Willmore-type flows, Calabi-type flows, the $Q$-curvature flow and the biharmonic
  map heat flow;
\item arbitrary even order: polyharmonic flows.
\end{itemize}
Almost every study of such a flow starts from the same statement: for a smooth initial datum at which
the equation is parabolic there is a unique smooth solution for a short time, and it depends
continuously on the datum. That statement is almost always justified by an appeal to ``standard
parabolic theory''.

For second-order equations and systems there is a large literature behind this appeal, for instance
\cite{Hamilton75,Eidelman,Solonnikov,Amann,Lunardi,Baker}. It is usually set in Euclidean domains, in
H\"older spaces or in an abstract semigroup framework, and the statement on a closed manifold has to be
extracted from it. For higher-order equations the situation is thinner.
\begin{itemize}[nosep]
\item Polden \cite{Polden} proposed a proof based on linear theory in parabolic Sobolev spaces and on
  the inverse function theorem.
\item Sharples \cite{Sharples} pointed out a gap in the convergence of the frozen-coefficient linear
  problems used there.
\item Mantegazza and Martinazzi \cite{MM} closed the gap for scalar quasilinear equations whose
  leading coefficient has the product form $E_1\otimes\cdots\otimes E_p$, with Polden's linear theory
  used as a black box. They also remarked that systems can be treated in the same way.
\end{itemize}

The purpose of this paper is to give one self-contained statement and proof that covers all of the
following at once:
\begin{itemize}[nosep]
\item equations of any even order $2p$;
\item sections of an arbitrary vector bundle with metric and compatible connection;
\item fully nonlinear dependence on all derivatives up to order $2p$;
\item parabolicity assumed only along the initial datum;
\item uniqueness, and smooth dependence on the initial datum.
\end{itemize}

\subsection{The main result}

In the informal statement below, $P_{u_0}=\partial F/\partial(\nabla^{2p}u)$ denotes the coefficient
of the top-order derivative in the linearization of $F$ at the initial datum.

\begin{theorem}[informal version of Theorems~\ref{thm:A}--\ref{thm:C}]\label{thm:informal}
Suppose that at every point of $M$ the coefficient $P_{u_0}$ satisfies the Legendre condition of
Definition~\ref{def:legendre}. Then:
\begin{enumerate}[nosep]
\item equation~\eqref{eq:intro} has a smooth solution on $M\times[0,T]$ for some $T>0$;
\item two smooth solutions along which the equation stays parabolic coincide;
\item the solution depends smoothly on $u_0$.
\end{enumerate}
\end{theorem}

The Legendre condition asks that the quadratic form
$\xi\mapsto(-1)^{p+1}\sum_{I,J}\ip{\xi_I}{P^{IJ}\xi_J}$ be positive definite on $V$-valued
$p$-tensors $\xi$. It is the higher-order analogue of the Legendre condition of the calculus of
variations.
\begin{itemize}[nosep]
\item It holds for $-(\Lap)^p$, and for every leading coefficient of product form, in particular for
  the class treated in \cite{MM} (Lemma~\ref{lem:legendre-examples}).
\item It is stable under the usual rearrangements of covariant derivatives
  (Remark~\ref{rem:rewrite}).
\item It is strictly stronger than Petrovskii's strong parabolicity in general, but equivalent to it
  for scalar second-order equations and in several other cases (Remark~\ref{rem:sos}).
\end{itemize}

\subsection{Strategy}

We keep the central idea of \cite{MM}. The nonlinear operator is a $C^1$ map between parabolic Sobolev
spaces; we linearize it at a smooth \emph{approximate} solution and apply the inverse function theorem.
Around this core we make five changes. Each one either widens the scope or removes a technical
ingredient.

\begin{enumerate}[label=(\arabic*)]
\item \emph{Approximate solutions from the formal power series.} We use the Taylor polynomial in $t$
  of the formal power series solution (Lemma~\ref{lem:formal}), rather than the solution of an
  auxiliary linear problem. This construction works verbatim for fully nonlinear equations.
\item \emph{No trace spaces.} The unknown is sought in the space $\Pp^m_0$ of functions vanishing at
  $t=0$, as a perturbation of the approximate solution. The error of the approximate solution is cut
  off near $t=0$ by a smooth cutoff (Lemma~\ref{lem:cutoff}). The sharp trace space of the parabolic
  Sobolev space never appears.
\item \emph{Linear theory proved, not quoted.} We obtain energy estimates at every integer level by
  differentiating the equation. Differentiation moves the problem from $V$ to $T^*M\otimes V$, which is
  why working with general bundles is natural even for scalar equations. The model operator
  $-(\Lap)^p$ is inverted by eigenfunction expansion, and general operators are reached by the method
  of continuity (Section~\ref{sec:linear}). Negative-order norms are defined by duality, and all
  orders are integers.
\item \emph{Regularity by Angenent's trick.} Pulling a solution back by the flows of horizontal lifts
  of vector fields produces a family of equations depending smoothly on a parameter. The implicit
  function theorem, applied at a smooth point, shows that the solution depends smoothly on the
  parameter, and this dependence is exactly spatial regularity (Section~\ref{sec:regularity}; compare
  \cite{Angenent90}). Consequently neither tame estimates nor linear theory with non-smooth
  coefficients is ever needed.
\item \emph{Uniqueness by linearization.} The difference of two solutions solves a linear equation
  with smooth coefficients, and uniqueness follows from the basic energy estimate
  (Section~\ref{sec:uniqueness}).
\end{enumerate}

\subsection{A proof designed to be formalized}

This paper is written as the blueprint of a formalization in Lean~4, on top of Mathlib \cite{mathlib}
and an existing library of differential geometry \cite{DGlib}. Several of the choices above were made
for that reason.

\begin{itemize}[nosep]
\item Every function space has integer order and is defined through covariant derivatives and time
  derivatives.
\item There are no interpolation spaces, trace theorems, Fourier multipliers or pseudodifferential
  operators.
\item Because the data are smooth, we may work in spaces of very high order and use only crude
  embeddings.
\item The linear theory is only ever applied to operators with smooth coefficients.
\item The ellipticity hypothesis is a pointwise algebraic condition that can be checked directly.
\item The final theorems are stated purely in terms of smooth sections, covariant derivatives and
  time derivatives. The Sobolev machinery is internal to the proofs.
\end{itemize}

\subsection{What is not covered}

The following are deliberately outside the scope of this paper.
\begin{itemize}[nosep]
\item Manifolds with boundary.
\item Noncompact manifolds.
\item Rough initial data and optimal regularity.
\item Strongly parabolic operators that admit no Legendre representation (Remark~\ref{rem:sos}).
\item Flows that are degenerate because of a gauge symmetry: the Ricci flow, the mean curvature and
  Willmore flows of immersions, the Yang--Mills flow. Our theorems apply to these after gauge fixing
  (for instance the Ricci--DeTurck flow, Section~\ref{ex:rdt}). The step that recovers the original
  flow from the gauge-fixed one is not treated here; see Section~\ref{sec:scope}.
\end{itemize}

\subsection{Organization}

\begin{itemize}[nosep]
\item Section~\ref{sec:setting}: the setting, the Legendre condition and the main theorems.
\item Section~\ref{sec:spaces}: the function spaces.
\item Section~\ref{sec:linear}: linear theory.
\item Section~\ref{sec:nonlinear}: the nonlinear theorems.
\item Section~\ref{sec:examples}: examples.
\item Section~\ref{sec:scope}: remarks on the scope.
\item Section~\ref{sec:formal}: the dependency structure and the imported facts.
\item Appendix~\ref{app:abstract}: the abstract tools.
\end{itemize}

% ===========================================================================
\section{Setting and statements}\label{sec:setting}
% ===========================================================================

\subsection{Bundles and covariant derivatives}

Throughout, $(M,g)$ is a closed Riemannian manifold, that is, compact and without boundary, of
dimension $n\ge 1$, with Riemannian measure $\mu$.

A \emph{Riemannian bundle} is a real vector bundle $W\to M$ with a fibre metric $\ip{\cdot}{\cdot}$ and
a metric connection $\nabla$. The Levi-Civita connection of $g$ and the connection of $W$ induce metric
connections on all bundles $\otimes^k T^*M\otimes W$. All of them are denoted by $\nabla$.

\paragraph{Iterated derivatives} For $u\in\Gamma(W)$ we define
$\nabla^k u\in\Gamma(\otimes^kT^*M\otimes W)$ inductively by $\nabla^0u=u$ and
\[
  (\nabla^k u)(X_1,\dots,X_k)=\bigl(\nabla_{X_1}(\nabla^{k-1}u)\bigr)(X_2,\dots,X_k).
\]
Thus the first slot records the derivative taken last, and $\nabla^k(\nabla^l u)=\nabla^{k+l}u$ as
tensors.

\paragraph{Frames and multi-indices} In a local orthonormal frame $(e_i)$ with dual coframe $(e^i)$, we
write $T_{i_1\dots i_k}\in W_x$ for the components of a $W$-valued $k$-tensor. For a multi-index
$I=(i_1,\dots,i_p)$ we write $e^I=e^{i_1}\otimes\cdots\otimes e^{i_p}$.

\paragraph{Permutations} A permutation $\sigma\in S_k$ acts on $W$-valued $k$-tensors by
$(\sigma T)(X_1,\dots,X_k)=T(X_{\sigma(1)},\dots,X_{\sigma(k)})$. This is a parallel bundle isometry.

\paragraph{Integration} We write $\ip{u}{v}=\int_M\ip{u}{v}\dmu$ and $\norm{u}$ for the $L^2$ inner
product and norm. The formal adjoint of $\nabla\colon\Gamma(W)\to\Gamma(T^*M\otimes W)$ is
\[
  \nabla^*\eta=-\sum_i(\nabla\eta)_{ii},
\]
where the derivative slot is contracted with the first slot of $\eta$. By the divergence theorem,
$\ip{\nabla u}{\eta}=\ip{u}{\nabla^*\eta}$. The connection Laplacian is
$\Lap u=-\sum_i(\nabla^2u)_{ii}$.

\paragraph{Time-dependent sections} A time-dependent section on $[0,T]$ is a section of the pullback
of $W$ to $M\times[0,T]$, and $C^\infty(M\times[0,T];W)$ denotes the smooth ones, smooth up to $t=0$
and $t=T$. The operator $\nabla$ acts in the $M$-variables. Since $g$ and the connection do not depend
on $t$, $\partial_t$ commutes with $\nabla$ and with all the bundle operations above.

\paragraph{Differential operators} For Riemannian bundles $W,W'$ and $r\in\N$, the space
$\Diff^r(W,W')$ consists of the operators
\[
  (Au)=\sum_{k=0}^{r}A_k(\nabla^k u),\qquad
  A_k\in C^\infty\bigl(M\times[0,T];\Hom(\otimes^kT^*M\otimes W,W')\bigr).
\]
Their coefficients may depend on $t$. By the Leibniz rule,
$\Diff^r\circ\Diff^s\subseteq\Diff^{r+s}$.

\begin{lemma}[commuting covariant derivatives]\label{lem:ricci}
\lean{ShortTime.permute_iteratedCovGrad}
Let $W$ be a Riemannian bundle, $r\ge2$ and $\sigma\in S_r$. There is an operator
$Q_\sigma\in\Diff^{r-2}(W,\otimes^rT^*M\otimes W)$ such that
\[
  \sigma(\nabla^r u)=\nabla^r u+Q_\sigma u\qquad\text{for all }u\in\Gamma(W).
\]
Its coefficients are universal polynomial expressions in the curvatures of $g$ and of $W$ and in their
covariant derivatives of order at most $r-2$.
\end{lemma}

\begin{proof}
Permutations act by parallel bundle isometries, which commute with $\nabla$ and preserve
$\Diff^{r-2}$. Hence the set of $\sigma$ for which the lemma holds contains the identity and is closed
under composition with adjacent transpositions. It therefore suffices to treat an adjacent
transposition $\tau_i=(i\ i{+}1)$ with $1\le i\le r-1$.

Write $\nabla^r u=\nabla^{i-1}(\nabla^2S)$ with $S=\nabla^{r-i-1}u$, a section of
$W'=\otimes^{r-i-1}T^*M\otimes W$. The Ricci identity states that
$(\nabla^2S)(X,Y)-(\nabla^2S)(Y,X)=R^{W'}(X,Y)S$, where the curvature $R^{W'}$ of $W'$ is expressed
through $\mathrm{Rm}_g$ and $R^W$. The permutation $\tau_i$ exchanges exactly the two slots of
$\nabla^2S$, so
\[
  \tau_i\nabla^ru-\nabla^ru=\nabla^{i-1}(\mathcal{R}S),
\]
where $\mathcal{R}$ is a zeroth-order bundle map. By the Leibniz rule,
$\nabla^{i-1}(\mathcal{R}S)=\sum_{a+b=i-1}(\nabla^a\mathcal R)*\nabla^{b+r-i-1}u$, which is an
operator of order $(i-1)+(r-i-1)=r-2$.
\end{proof}

\subsection{Jets and nonlinear operators}

\begin{definition}[covariant jets]\label{def:jet}
\lean{ShortTime.CovJetFiber, ShortTime.covJet}
For a Riemannian bundle $W$ we set
\[
  \Jet^rW=\bigoplus_{k=0}^r\otimes^kT^*M\otimes W,\qquad
  j^ru=(u,\nabla u,\dots,\nabla^ru),
\]
and call $j^ru$ the \emph{covariant $r$-jet} of $u$. Elements of $\Jet^rW$ are written
$\zeta=(\zeta_0,\dots,\zeta_r)$, and $\Jet^rW$ carries the direct-sum fibre metric.
\end{definition}

\begin{definition}[equation data]\label{def:data}
\lean{ShortTime.EquationData, ShortTime.localDom, ShortTime.localRep}
\uses{def:jet}
The data of the equation are:
\begin{itemize}[nosep]
\item an integer $p\ge1$, a time $T_0>0$ and a Riemannian bundle $V$;
\item an open subset $\Dom\subseteq[0,T_0]\times\Jet^{2p}V$ (open in the relative topology);
\item a smooth map $F\colon\Dom\to V$ with $F(t,\zeta)\in V_x$ whenever $\zeta\in\Jet^{2p}_xV$.
\end{itemize}
\end{definition}

For a time-dependent section $u$ we say that $u$ \emph{has jets in} a set $\mathcal S$ if
$(t,j^{2p}u(x,t))\in\mathcal S$ for all $(x,t)$. When $u$ has jets in $\Dom$ we write
$F[u](x,t)=F\bigl(t,j^{2p}u(x,t)\bigr)$.

\begin{definition}\label{def:solution}
\lean{ShortTime.JointlySmoothOn, ShortTime.EquationData.IsSolution, ShortTime.EquationData.IsParabolicSolution}
\uses{def:data, def:jet}
Let $u_0\in C^\infty(M;V)$ and $T\in(0,T_0]$. A \emph{solution on $[0,T]$ with initial datum $u_0$}
is a section $u\in C^\infty(M\times[0,T];V)$ that has jets in $\Dom$ and satisfies
$\partial_tu=F[u]$ and $u(\cdot,0)=u_0$.
\end{definition}

\begin{definition}[linearization]\label{def:linearization}
\lean{ShortTime.EquationData.partialDeriv, ShortTime.EquationData.principal}
\uses{def:data}
For $0\le k\le 2p$ let
$D_kF(t,\zeta)\in\Hom(\otimes^kT^*_xM\otimes V_x,V_x)$ denote the derivative of $F$ along the summand
$\otimes^kT^*M\otimes V$ of the fibre. For a smooth $u$ with jets in $\Dom$ we define the linearization
\begin{equation}\label{eq:linearization}
  L_u v=\sum_{k=0}^{2p}D_kF\bigl(t,j^{2p}u\bigr)(\nabla^kv),
  \qquad\text{so that}\qquad
  \frac{d}{d\varepsilon}\Big|_{\varepsilon=0}F[u+\varepsilon v]=L_uv,
\end{equation}
and its principal coefficient $P_u=D_{2p}F(t,j^{2p}u)$.
\end{definition}

\subsection{The Legendre condition}

\begin{definition}[Legendre condition]\label{def:legendre}
\lean{ShortTime.legendreForm, ShortTime.LegendreAt, ShortTime.EquationData.IsParabolicAt}
\uses{def:linearization}
Let $W$ be a Riemannian bundle, $x\in M$ and
$P\in\Hom(\otimes^{2p}T^*_xM\otimes W_x,W_x)$.
\begin{enumerate}[nosep]
\item Define $\hatP\in\End(\otimes^pT^*_xM\otimes W_x)$ by
  \[
    \ip{\xi}{\hatP\eta}=\sum_{I,J}\ip{\xi_I}{P(e^I\otimes e^J\otimes\eta_J)},
  \]
  where $I$ and $J$ range over $\{1,\dots,n\}^p$ and $(e_i)$ is an orthonormal basis of $T_xM$. The
  right-hand side does not depend on the basis.
\item The \emph{Legendre modulus} of $P$ is
  \[
    \kappa(P)=\min_{\abs{\xi}=1}(-1)^{p+1}\ip{\xi}{\hatP\xi}.
  \]
\item $P$ satisfies the \emph{Legendre condition} if $\kappa(P)>0$.
\end{enumerate}
The \emph{parabolicity domain} of $F$ is
\[
  \Par=\{(t,\zeta)\in\Dom:\ \kappa\bigl(D_{2p}F(t,\zeta)\bigr)>0\}.
\]
A solution is \emph{parabolic} if it has jets in $\Par$.
\end{definition}

The sign $(-1)^{p+1}$ is chosen so that $\partial_tu=-(\Lap)^pu$, the ``heat equation of order
$2p$'', is parabolic; see Lemma~\ref{lem:legendre-examples}(b).

\begin{lemma}\label{lem:legendre-basic}
\lean{ShortTime.legendreForm_add, ShortTime.legendreForm_smul, ShortTime.LegendreAt.combo, ShortTime.isOpen_parabolicDomain, ShortTime.legendre_tube}
\uses{def:legendre}
\begin{enumerate}[label=(\roman*),nosep]
\item The map $P\mapsto\hatP$ is linear. Moreover $\kappa$ is superadditive and positively
  homogeneous: $\kappa(P+P')\ge\kappa(P)+\kappa(P')$ and $\kappa(cP)=c\,\kappa(P)$ for $c\ge0$. In
  particular each set $\{\kappa\ge\lambda\}$ is convex.
\item The set $\Par$ is open in $[0,T_0]\times\Jet^{2p}V$.
\item Let $K\subseteq\Par$ be compact. There are $\delta>0$ and $\kappa_K>0$ with the following
  property. Suppose $(t',\zeta')\in[0,T_0]\times\Jet^{2p}V$ and there is $(t,\zeta)\in K$ over the same
  point of $M$ with $\abs{t'-t}\le\delta$ and $\abs{\zeta'-\zeta}\le\delta$. Then
  $(t',\zeta')\in\Par$ and $\kappa(D_{2p}F(t',\zeta'))\ge\kappa_K$.
\end{enumerate}
\end{lemma}

\begin{proof}
(i) The map $P\mapsto\hatP$ is linear by definition, and a minimum of a sum is at least the sum of the
minima.

(ii) The function $(t,\zeta)\mapsto\kappa(D_{2p}F(t,\zeta))$ is continuous on $\Dom$, because $\kappa$
is continuous in $P$ and in the base point, and $\Par$ is where it is positive.

(iii) Let $K_\delta$ be the set of points $(t',\zeta')$ described in the statement. Then $K_\delta$ is
compact, decreases with $\delta$, and $\bigcap_\delta K_\delta=K$. The continuous function
$\kappa\circ D_{2p}F$ is at least $2\kappa_K:=\min_K\kappa\circ D_{2p}F>0$ on $K$, hence at least
$\kappa_K$ on an open neighbourhood $U\subseteq\Dom$ of $K$. Since the $K_\delta$ are compact, they are
contained in $U$ for $\delta$ small.
\end{proof}

\begin{lemma}[examples]\label{lem:legendre-examples}
\lean{ShortTime.legendre_product, ShortTime.isDiffOp_lapIter_sub_modelTop, ShortTime.legendreAt_modelTop, ShortTime.polyTrace, ShortTime.legendreAt_polyTrace, ShortTime.productCLM, ShortTime.mmPairing}
\uses{def:legendre, lem:ricci}
\begin{enumerate}[label=(\alph*),nosep]
\item \emph{Product structure.} Let $E_1,\dots,E_p$ be symmetric forms on $T^*_xM$ with
  $E_\ell\ge e_\ell\,g^{-1}$ for constants $e_\ell>0$. Let $B\in\End(W_x)$ be symmetric with
  $B\ge b>0$. Let
  \[
    P(X)=(-1)^{p+1}\sum E_1^{i_1j_1}\cdots E_p^{i_pj_p}\,B\,X_{i_1\dots i_pj_1\dots j_p}.
  \]
  Then $\kappa(P)\ge b\,e_1\cdots e_p$.
\item \emph{Polyharmonic operator.} There is $Q_0\in\Diff^{2p-2}(W,W)$ with
  \[
    -(\Lap)^pu=P_0(\nabla^{2p}u)+Q_0u,\qquad
    P_0(X)=(-1)^{p+1}\sum_{I}X_{II},
  \]
  and $\kappa(P_0)=1$.
\end{enumerate}
\end{lemma}

\begin{proof}
(a) Identify $E_\ell$ with a symmetric endomorphism of $T^*_xM$ via $g$. Then
$(-1)^{p+1}\hatP=E_1\otimes\cdots\otimes E_p\otimes B$ on $\otimes^pT^*_xM\otimes W_x$. A tensor
product of positive symmetric operators is positive, and its spectrum consists of products of
eigenvalues of the factors.

(b) Since contraction commutes with $\nabla$,
\[
  (\Lap)^pu=(-1)^p\sum_{i_1,\dots,i_p}(\nabla^{2p}u)_{i_1i_1i_2i_2\dots i_pi_p}.
\]
Choose $\sigma\in S_{2p}$ that moves the slot pairs $(1,2),(3,4),\dots$ to the pairs
$(1,p{+}1),(2,p{+}2),\dots$. Lemma~\ref{lem:ricci} applied to $\sigma$ yields
$-(\Lap)^pu=(-1)^{p+1}\sum_I(\nabla^{2p}u)_{II}+Q_0u$. Finally, $P_0$ is the case $E_\ell=g^{-1}$,
$B=\id$ of (a), and there $(-1)^{p+1}\widehat{P_0}=\id$.
\end{proof}

\begin{remark}[rewriting the top-order term]\label{rem:rewrite}
The Legendre condition depends on how the top-order derivatives are arranged in $F$, not only on the
equation. By Lemma~\ref{lem:ricci} this dependence is harmless. Let $P'$ be any smooth coefficient
whose restriction to symmetric tensors agrees with that of $D_{2p}F$. Then
$(P'-D_{2p}F)(\nabla^{2p}u)$ equals a lower-order expression in $u$, because
$\nabla^{2p}u-\operatorname{Sym}\nabla^{2p}u=-\frac{1}{(2p)!}\sum_\sigma Q_\sigma u$. So $F$ can be
replaced by a function $F'$ with $F'[u]=F[u]$ for every $u$ and $D_{2p}F'=P'$.

Our theorems therefore apply whenever \emph{some} such rearrangement satisfies the Legendre
condition. For example, the operators of \cite{MM} are written with the pairing
$A^{i_1j_1\cdots i_pj_p}\nabla^{2p}_{i_1j_1\cdots i_pj_p}$. After moving to the split ordering of
Lemma~\ref{lem:legendre-examples}(a) they satisfy the Legendre condition.
\end{remark}

\begin{remark}[Legendre condition and strong parabolicity]\label{rem:sos}
\emph{Legendre implies strong parabolicity.} Taking $\xi=\eta^{\otimes p}\otimes v$ in
Definition~\ref{def:legendre} shows that the Legendre condition implies strong (Petrovskii)
parabolicity:
\[
  (-1)^{p+1}\ip{v}{P(\eta^{\otimes 2p}\otimes v)}\ge\kappa(P)\abs{\eta}^{2p}\abs{v}^2 .
\]

\emph{Characterization for scalar equations.} In the scalar case $W=\R$, and in view of
Remark~\ref{rem:rewrite}, an operator admits a Legendre rearrangement if and only if
$(-1)^{p+1}\sigma_P(\eta)-\lambda\abs\eta^{2p}$ is a sum of squares of forms of degree $p$ in $\eta$
for some $\lambda>0$. Here $\sigma_P(\eta)=P(\eta^{\otimes2p})$ is the symbol.

\emph{Consequences of Hilbert's theorem.} By Hilbert's theorem \cite{Hilbert}, every nonnegative form
is a sum of squares exactly for binary forms, quadratic forms, and ternary quartics. Hence for scalar
equations with $p=1$, or with $n\le2$, or with $n=3$ and $p=2$, strong parabolicity and the existence
of a Legendre rearrangement are equivalent. In general they are not: a small perturbation of the
Motzkin form \cite{Motzkin} gives a strongly parabolic scalar operator of order six in dimension
three with no Legendre rearrangement. Such operators are outside the scope of this paper.
\todo{double-check the precise form of this remark, in particular strict positivity versus SOS.}
\end{remark}

\subsection{How to check the hypothesis in practice}\label{sec:howto}

General statements about ellipticity are easy to get lost in, so here is what the Legendre condition
amounts to and how to verify it for a given equation.

\paragraph{The underlying idea} Freeze the coefficients and insert a plane wave
$u=e^{i\ip{\xi}{x}}a(t)$. Each covariant derivative becomes a factor of $\xi$, and the top-order term
acts on $a$ like a matrix of size $\abs\xi^{2p}$, the \emph{symbol}. Parabolicity means that
high-frequency waves are damped at a rate comparable to $\abs\xi^{2p}$. For example:
\begin{itemize}[nosep]
\item $\partial_tu=\Delta u$ has symbol $-\abs\xi^2$, so high frequencies are damped;
\item $\partial_tu=+\Delta^2u$ has symbol $+\abs\xi^4$, so high frequencies grow: this is a backward
  heat equation.
\end{itemize}
There are three standard ways to make this precise.
\begin{enumerate}[label=(\arabic*),nosep]
\item \emph{Strong parabolicity}: the symbol is uniformly negative. This only tests pure powers
  $\xi\otimes\cdots\otimes\xi$.
\item \emph{G\aa rding's inequality}: $-\ip{\psi}{L\psi}\ge c\norm\psi_{H^p}^2-C\norm\psi^2$, that is,
  the operator dissipates energy up to lower-order errors. This is the only form in which ellipticity
  enters the proofs below: everything after Lemma~\ref{lem:garding} uses nothing else.
\item \emph{The Legendre condition}: the matrix $(-1)^{p+1}\hatP$ is positive definite on
  \emph{all} $p$-tensors, not only on pure powers.
\end{enumerate}

\paragraph{Why Legendre is the convenient version} Integrating
$\int\ip{\psi}{P(\nabla^{2p}\psi)}$ by parts $p$ times moves half of the derivatives onto the other
factor. The result is $\int\ip{\nabla^p\psi}{\hatP\,\nabla^p\psi}$ up to lower-order terms
(Lemma~\ref{lem:ibp}). Since $\nabla^p\psi$ can be an essentially arbitrary $p$-tensor, positivity of
$\hatP$ on all $p$-tensors is exactly what this integration by parts needs.

Deriving G\aa rding's inequality from strong parabolicity alone is also possible, but it requires
Fourier analysis, freezing of coefficients and a partition of unity. In Fourier variables
$\nabla^p\psi$ becomes a pure power, but only for constant coefficients. The difference between the
two conditions is precisely the difference between ``positive on pure powers'' and ``positive on
everything''.

\paragraph{A recipe}
\begin{enumerate}[label=(\arabic*),nosep]
\item Linearize the equation and take the coefficient $P$ of $\nabla^{2p}u$.
\item Arrange its $2p$ indices so that the first $p$ slots are paired with the last $p$ slots,
  commuting covariant derivatives where necessary. This changes the equation only by lower-order
  terms (Remark~\ref{rem:rewrite}).
\item Check that the resulting matrix, multiplied by $(-1)^{p+1}$, is positive definite on
  $V$-valued $p$-tensors.
\end{enumerate}
For the top-order terms occurring in geometric flows, namely $-(-\Delta)^p$ or a product
$E_1\otimes\cdots\otimes E_p$ of positive definite forms, step~(3) is automatic by
Lemma~\ref{lem:legendre-examples}.

\begin{example}[the arrangement matters]\label{ex:bilaplacian}
Consider $\partial_tu=-\Delta^2u$ for scalar $u$, so $p=2$ and $(-1)^{p+1}=-1$.
\begin{itemize}[nosep]
\item \emph{Adjacent pairing.} Writing
  $-\Delta^2u=-g^{ab}g^{cd}\nabla_a\nabla_b\nabla_c\nabla_du$ and splitting the slots as $(a,b)$ and
  $(c,d)$ gives the form $\xi\mapsto(\tr_g\xi)^2$ on $2$-tensors. It vanishes on trace-free $\xi$, so
  the Legendre condition fails for this way of writing the operator.
\item \emph{Split pairing.} Exchanging the two middle derivatives costs only a curvature term of
  order two (Lemma~\ref{lem:ricci}). It gives
  $-\Delta^2u=-g^{ac}g^{bd}\nabla_a\nabla_b\nabla_c\nabla_du+(\text{lower order})$. The form is now
  $\xi\mapsto g^{ac}g^{bd}\xi_{ab}\xi_{cd}=\abs\xi^2$, so the Legendre condition holds with
  $\kappa=1$.
\end{itemize}
This is Lemma~\ref{lem:legendre-examples}(b) for $p=2$.
\end{example}

If an equation is strongly parabolic but admits no Legendre arrangement (Remark~\ref{rem:sos}), one
can replace Lemma~\ref{lem:garding} by the general G\aa rding inequality. Nothing downstream changes.

\subsection{Main results}

Set
\[
  m_0=\Bigl\lfloor\frac{n}{2p}\Bigr\rfloor+4 .
\]
The spaces $H^{k}(M;V)$ and $\Pp^m(V)$ used in Theorem~\ref{thm:C} are defined in
Section~\ref{sec:spaces}.

\begin{mainthm}[existence and regularity]\label{thm:A}
\lean{ShortTime.exists_solution}
\uses{def:solution, def:legendre}
Let $u_0\in C^\infty(M;V)$ with $(0,j^{2p}u_0(x))\in\Par$ for every $x\in M$. Then there are
$T\in(0,T_0]$ and a parabolic solution $u\in C^\infty(M\times[0,T];V)$ with initial datum $u_0$.
\end{mainthm}

\begin{mainthm}[uniqueness]\label{thm:B}
\lean{ShortTime.solution_unique, ShortTime.EquationData.solution_unique_cov}
\uses{def:solution, def:legendre}
Two parabolic solutions on $[0,T]$ with the same initial datum coincide.
\end{mainthm}

\begin{mainthm}[smooth dependence]\label{thm:C}
\lean{ShortTime.solution_smooth_dependence}
\uses{def:solution, def:legendre, def:sobolev, def:parabolic}
Let $u$ be a parabolic solution on $[0,T]$ with initial datum $u_0$, and let $m\ge m_0$. There are
$\delta>0$ and a $C^\infty$ map
\[
  \mathcal S\colon\{a\in H^{2pm}(M;V):\norm{a-u_0}_{H^{2pm}}<\delta\}\to\Pp^m(V)
\]
with $\mathcal S(u_0)=u$ and the following property: for every smooth $a$ in the domain,
$\mathcal S(a)$ is a parabolic solution on $[0,T]$ with initial datum $a$.
\end{mainthm}

\begin{corollary}[continuous dependence, sequential form]\label{cor:C}
\lean{ShortTime.solution_continuous_dependence, ShortTime.solution_continuous_dependence', ShortTime.solution_continuous_dependence_of_covariant, ShortTime.EquationData.continuous_dependence_cov, ShortTime.EquationData.traceRegularity, ShortTime.EquationData.smoothRep_time_trace}
\uses{def:solution, def:legendre}
Let $u$ be a parabolic solution on $[0,T]$ with initial datum $u_0$. If $u_0^{(k)}\to u_0$ in
$C^\infty(M;V)$, then for all large $k$ the parabolic solution $u^{(k)}$ with initial datum
$u_0^{(k)}$ exists on $[0,T]$, and $u^{(k)}\to u$ in $C^\infty(M\times[0,T];V)$.
\end{corollary}

The formal version of Corollary~\ref{cor:C} states the convergence of every covariant derivative
$\nabla^j(u^{(k)}-u)$, uniformly on $M\times[0,T]$. Convergence of mixed derivatives
$\nabla^j\partial_t^i(u^{(k)}-u)$ then follows by induction on $i$ from the equation, since the jets
of $u$ lie in a compact subset of $\mathcal D$. Theorem~\ref{thm:C} itself, which involves the
spaces $H^{2pm}$ and $\Pp^m$, is not among the trusted formal statements
(Section~\ref{sec:formal}); it is used in the proof of Corollary~\ref{cor:C}.

\begin{corollary}[maximal solutions]\label{cor:maximal}
\lean{ShortTime.exists_maximal_solution, ShortTime.EquationData.isParabolicSolution_of_le}
\uses{def:solution}
Under the hypothesis of Theorem~\ref{thm:A}, let $T^*$ be the supremum of the times $T\le T_0$ for
which a parabolic solution on $[0,T]$ exists. There is a unique section $u$ on $M\times[0,T^*)$, or on
$M\times[0,T_0]$ if $T^*=T_0$ is attained, whose restriction to every $[0,T]$ with $T<T^*$ is a
parabolic solution.
\end{corollary}

\begin{proof}
\uses{thm:A, thm:B}
By Theorem~\ref{thm:B}, parabolic solutions on different intervals agree where both are defined.
\end{proof}

\begin{corollary}[Mantegazza--Martinazzi {\cite[Thm.~1.1]{MM}}]\label{cor:MM}
\lean{ShortTime.mantegazza_martinazzi, ShortTime.MMData, ShortTime.MMData.F, ShortTime.MMData.toEquationData, ShortTime.MMData.rearranged, ShortTime.MMData.rearranged_isSolution_iff, ShortTime.MMData.rearranged_isParabolicAt, ShortTime.MMData.isParabolicSolution_rearranged, ShortTime.exists_ricci_jetMap, ShortTime.mmPairing_eq_productCLM_add}
\uses{def:solution}
Let $V=M\times\R$ with its flat connection, and let
\[
  F(t,\zeta)=A(t,\zeta_0,\dots,\zeta_{2p-1})(\zeta_{2p})+b(t,\zeta_0,\dots,\zeta_{2p-1})
\]
on $\Dom=[0,T_0]\times\Jet^{2p}V$. Suppose $A$ has the product structure
$A=(-1)^{p-1}E_1\otimes\cdots\otimes E_p$ of \cite{MM}, and each $E_\ell$ is locally uniformly
positive definite, meaning that on every compact subset of its domain it is bounded below by a
positive multiple of $g^{-1}$. After the rearrangement of Remark~\ref{rem:rewrite} we have
$\Par=\Dom$. In particular, every smooth initial datum has a unique smooth solution for short time,
and it depends continuously on the initial datum in $C^\infty$.
\end{corollary}

\begin{proof}
\uses{thm:A, thm:B, cor:C, lem:legendre-examples, lem:ricci}
By Remark~\ref{rem:rewrite} (made exact by the Ricci identity, Lemma~\ref{lem:ricci}: permuting the
$2p$ derivative slots changes $\nabla^{2p}u$ by a smooth function of the jet of order $2p-2$) we
replace $F$ by $F'$ with the same value on jets of sections, whose principal coefficient is the split
pairing $(-1)^{p+1}E_1\otimes\cdots\otimes E_p$. By Lemma~\ref{lem:legendre-examples}(a), with the
sign $(-1)^{p-1}=(-1)^{p+1}$, every point of $\Dom$ lies in $\Par$ for $F'$. Solutions of $F$ and
$F'$ coincide, so Theorems~\ref{thm:A}, \ref{thm:B} and Corollary~\ref{cor:C} apply to $F$. In the
formalization the $E_\ell$ are only assumed positive definite at each point; locally uniform
positivity then follows by continuity.
\end{proof}

% ===========================================================================
\section{Function spaces}\label{sec:spaces}
% ===========================================================================

In this section $W$ is a Riemannian bundle over $M$. Constants $C$ may change from line to line. They
depend only on the quantities indicated and on $(M,g)$ and $W$.

\subsection{Sobolev spaces of integer order}

\begin{definition}\label{def:sobolev}
\lean{ShortTime.sobolevNormSq, ShortTime.negSobolevNorm, ShortTime.hNormSq}
Let $k\in\N$ and $u,f\in C^\infty(M;W)$.
\begin{enumerate}[nosep]
\item The positive-order norm is $\norm{u}_{H^k}^2=\sum_{i=0}^k\norm{\nabla^iu}^2$. The space
  $H^k(M;W)$ is the completion of $C^\infty(M;W)$ in this norm; it embeds injectively in $L^2$ because
  $\nabla$ is closable.
\item The negative-order norm is
  \[
    \norm{f}_{H^{-k}}=\sup\{\ip{f}{\varphi}:\varphi\in C^\infty(M;W),\ \norm{\varphi}_{H^k}\le1\}.
  \]
\end{enumerate}
We only use these norms on smooth sections.
\end{definition}

\begin{lemma}\label{lem:sobolev-basic}
\lean{ShortTime.sobolevNormSq_succ, ShortTime.inner_le_negSobolevNorm, ShortTime.exists_negSobolevNorm_covGrad_le, ShortTime.exists_sobolevNormSq_opField_le, ShortTime.exists_negSobolevNorm_opField_covGrad_le, ShortTime.sobolevNorm_diffOp_le}
\uses{def:sobolev}
For smooth sections:
\begin{enumerate}[label=(\roman*),nosep]
\item $\norm{u}_{H^{k+1}}^2=\norm{u}^2+\norm{\nabla u}_{H^k}^2$ for $k\ge0$.
\item $\abs{\ip{f}{\varphi}}\le\norm f_{H^{-k}}\norm\varphi_{H^k}$. Moreover $s\mapsto\norm f_{H^s}$
  is nondecreasing on $\mathbb Z$.
\item For every $s\in\mathbb Z$ there is $C$ with $\norm{\nabla f}_{H^{s-1}}\le C\norm f_{H^s}$.
\item Let $A\in\Diff^r(W,W')$ be a time-independent operator, or a single time slice of a
  time-dependent one. If $s\ge0$ then $\norm{Au}_{H^s}\le C\norm{u}_{H^{s+r}}$. If $r=0$ this also
  holds for $s<0$. The constant depends on bounds for finitely many derivatives of the coefficients of
  $A$.
\end{enumerate}
\end{lemma}

\begin{proof}
(i) follows from the definition, since $\nabla^{i+1}u=\nabla^i(\nabla u)$. The inequality in (ii) is
the definition, and monotonicity follows from $\norm\varphi_{H^{k}}\le\norm\varphi_{H^{k+1}}$.

(iii) For $s\ge1$ this is (i). For $s=-k\le0$, observe that
$\ip{\nabla f}{\varphi}=\ip{f}{\nabla^*\varphi}$ and
$\norm{\nabla^*\varphi}_{H^k}\le\sqrt n\,\norm{\nabla\varphi}_{H^k}\le\sqrt n\,\norm{\varphi}_{H^{k+1}}$,
because $\nabla^*$ is a contraction of $\nabla$ and contractions commute with $\nabla$.

(iv) For $s\ge0$ this follows from the Leibniz rule. For $s<0$ and $r=0$ use duality:
$\ip{Au}{\varphi}=\ip{u}{A^\dagger\varphi}$, where $A^\dagger$ is the pointwise adjoint, and apply the
case $s\ge0$ to $A^\dagger$.
\end{proof}

\begin{lemma}[interpolation]\label{lem:interp}
\lean{ShortTime.jetNorm_le_interp, ShortTime.exists_jetSum_pred_le, ShortTime.exists_sobolevNormSq_pred_le, ShortTime.exists_sobolevNormSq_le}
\uses{def:sobolev}
For integers $0<j<k$ and every $\varepsilon>0$ there is $C=C(n,j,k,\varepsilon)$ such that
\[
  \norm{\nabla^ju}\le\varepsilon\norm{\nabla^ku}+C\norm u\qquad\text{for all smooth }u.
\]
Consequently, for $k\ge1$ and with $N_k(u)=\sum_{i\le k}\norm{\nabla^iu}$,
\[
  N_{k-1}(u)\le\varepsilon\norm{\nabla^ku}+C\norm u,\qquad
  \norm{u}_{H^{k-1}}^2\le\varepsilon\norm{u}_{H^k}^2+C\norm u^2,\qquad
  \norm{u}_{H^k}^2\le2\norm{\nabla^ku}^2+C\norm u^2.
\]
\end{lemma}

\begin{proof}
Let $a_i=\norm{\nabla^iu}$. Integrating by parts,
\[
  a_i^2=\ip{\nabla(\nabla^{i-1}u)}{\nabla^iu}=\ip{\nabla^{i-1}u}{\nabla^*\nabla^iu}\le\sqrt n\,a_{i-1}a_{i+1}.
\]
If some $a_i$ with $i\ge1$ vanishes, then $a_j=0$ for all $j\ge1$ and there is nothing to prove.
Otherwise $c_i=\log a_i+\tfrac14(\log n)\,i^2$ is a convex sequence. Hence
$a_j\le n^{j(k-j)/4}a_0^{1-j/k}a_k^{j/k}$, and Young's inequality gives the claim. The consequences
follow by summing over $j$, using $\norm u_{H^{k-1}}^2\le N_{k-1}(u)^2$ and
$\norm{\nabla^ku}\le\norm u_{H^k}$.
\end{proof}

Note that $\norm u_{H^k}\le N_k(u)\le\sqrt{k+1}\,\norm u_{H^k}$. We set $N_k=N_0$ for $k<0$.

\begin{lemma}[permuting derivatives in $L^2$]\label{lem:ricci-l2}
\lean{ShortTime.exists_perm_bound, ShortTime.exists_swap_defect, ShortTime.toModel_iteratedCovGrad_domDomCongr}
\uses{def:sobolev}
Let $r\in\N$ and let $\sigma\in S_r$ act on the $r$ derivative slots of $\nabla^ru$. There is $C$ with
\[
  \norm{\sigma(\nabla^ru)-\nabla^ru}\le C\,N_{r-2}(u)\qquad\text{for all smooth }u.
\]
\end{lemma}

\begin{proof}
Permutations of slots are fibrewise isometries. Hence the set of $\sigma$ for which the estimate
holds contains the identity and is closed under products, since
$\norm{\sigma\sigma'\nabla^ru-\nabla^ru}\le\norm{\sigma'\nabla^ru-\nabla^ru}+\norm{\sigma\nabla^ru-\nabla^ru}$.
It therefore suffices to treat an adjacent transposition $\tau_i$. As in the proof of
Lemma~\ref{lem:ricci}, $\tau_i\nabla^ru-\nabla^ru=\nabla^{i-1}(\mathcal RS)$ with
$S=\nabla^{r-i-1}u$ and $\mathcal R$ of order zero: the Ricci identity handles the two slots created
by $\nabla^2S$, and $\nabla$ commutes with permutations of the slots it does not create. By the
Leibniz rule $\abs{\nabla^{i-1}(\mathcal RS)}\le C\sum_{b\le i-1}\abs{\nabla^bS}$ pointwise, and
integrating gives the claim.
\end{proof}

\begin{fact}[Sobolev embedding {\cite{Aubin,Hebey}}]\label{fact:sobolev}
\uses{def:sobolev}
\source{DG: Analysis/Sobolev/Embedding/Tensor (spectral norms; needs norm equivalence)}
If $k>n/2+i$, then $\sup_M\abs{\nabla^iu}\le C\norm{u}_{H^k}$ for smooth $u$. Hence elements of $H^k$
are $C^i$ sections.
\end{fact}

\begin{fact}[weak and strong derivatives]\label{fact:weak}
\uses{def:sobolev}
\source{DG: Analysis/Sobolev/Intrinsic (to check)}
Let $Y_1,\dots,Y_q$ be smooth vector fields spanning $T_xM$ at every $x$. Let $u\in L^2(M;W)$, and
suppose that for every word $(a_1,\dots,a_d)$ with $d\le k$ the weak derivative
$\nabla_{Y_{a_1}}\cdots\nabla_{Y_{a_d}}u$ lies in $L^2$, defined by duality with the formal adjoint
$\nabla_Y^\dagger\psi=-\nabla_Y\psi-(\operatorname{div}Y)\psi$. Then $u\in H^k(M;W)$ and
\[
  \norm{u}_{H^k}\le C\sum_{d\le k}\ \sum_{a_1,\dots,a_d}\norm{\nabla_{Y_{a_1}}\cdots\nabla_{Y_{a_d}}u}.
\]
The same holds for time-dependent sections with $L^2(0,T;\cdot)$ norms. This follows from Friedrichs
mollification in charts.
\end{fact}

\subsection{Parabolic Sobolev spaces}

Fix $T\in(0,T_0]$ and $m\in\N$.

\begin{definition}\label{def:parabolic}
\lean{ShortTime.ParabolicSpace, ShortTime.ParabolicSpace.zeroAtStart}
\uses{def:sobolev}
For $u\in C^\infty(M\times[0,T];W)$ let
\[
  \norm{u}_{\Pp^m}^2=\sum_{\substack{j,i\ge0\\2pj+i\le2pm}}\int_0^T\norm{\partial_t^j\nabla^iu}^2dt
  =\sum_{j=0}^m\int_0^T\norm{\partial_t^ju}_{H^{2p(m-j)}}^2dt .
\]
\begin{itemize}[nosep]
\item $\Pp^m(W)$ is the completion of $C^\infty(M\times[0,T];W)$ in this norm. It embeds in
  $L^2(M\times[0,T];W)$.
\item $\Pp^m_0(W)$ is the closure in $\Pp^m(W)$ of $\{u\in C^\infty:\ u(\cdot,0)=0\}$.
\end{itemize}
Thus $\Pp^0=L^2$, and $\Pp^m_0$ is a closed subspace of $\Pp^m$, so it is a Hilbert space.
\end{definition}

\begin{lemma}\label{lem:parabolic-basic}
\lean{ShortTime.ParabolicSpace.timeDeriv, ShortTime.ParabolicSpace.jet, ShortTime.ParabolicSpace.constExt, ShortTime.ParabolicSpace.evalJet, ShortTime.ParabolicSpace.evalJet_ofSmooth, ShortTime.ParabolicSpace.continuousOn_evalJet_bundle, ShortTime.ParabolicSpace.evalJet_eq_zero_of_mem_zeroAtStart, ShortTime.ParabolicSpace.eq_zero_of_ofSmooth_eq_zero, ShortTime.ParabolicSpace.levelIncl, ShortTime.ParabolicSpace.evalJet_levelIncl, ShortTime.sobolevNormSq_le_integral, ShortTime.sobolevNormSq_timeDerivIter_le_parabolicInner, ShortTime.parabolicInner_iteratedCovGrad_le, ShortTime.parabolicInner_timeDeriv_le}
\uses{def:parabolic, lem:sobolev-basic}
\begin{enumerate}[label=(\roman*),nosep]
\item $\partial_t\colon\Pp^{m+1}\to\Pp^m$ and $\nabla^k\colon\Pp^m\to\Pp^{m-1}$ for $k\le2p$ are
  bounded. Hence $j^{2p}\colon\Pp^m(W)\to\Pp^{m-1}(\Jet^{2p}W)$ is bounded.
\item $\Pp^{m+1}\hookrightarrow\Pp^m$ with norm at most $1$.
\item Every $A\in\Diff^{2p}(W,W')$ whose coefficients are smooth on $M\times[0,T]$ is bounded
  $\Pp^m\to\Pp^{m-1}$. In particular $\partial_t-L$ is bounded $\Pp^m\to\Pp^{m-1}$ for
  $L\in\Diff^{2p}(W,W)$.
\item For $a\in H^{2pm}(M;W)$, the constant extension $(Ea)(x,t)=a(x)$ satisfies
  $\norm{Ea}_{\Pp^m}=\sqrt T\,\norm a_{H^{2pm}}$.
\item If $m\ge1$, then $u\mapsto u(\cdot,t)$ extends to a bounded map $\Pp^m\to L^2(M;W)$ that is
  continuous in $t$. Moreover $u(\cdot,0)=0$ for $u\in\Pp^m_0$.
\end{enumerate}
\end{lemma}

\begin{proof}
(i) and (ii) are immediate from the definition: for instance
$\partial_t^j\nabla^i\nabla^ku=\partial_t^j\nabla^{i+k}u$ and $2pj+i+k\le2pm$ when
$2pj+i\le2p(m-1)$ and $k\le2p$. (iii) follows from the Leibniz rule in $x$ and $t$. (iv) is a direct
computation.

(v) For a Hilbert space $X$ and $h\in C^1([0,T];X)$ we have
$\sup_t\norm{h(t)}_X\le C_T\bigl(\norm h_{L^2X}+\norm{h'}_{L^2X}\bigr)$. Apply this with $X=L^2$ and
$h=u$.
\end{proof}

\begin{lemma}[embedding]\label{lem:embedding}
\lean{ShortTime.exists_fiberNormSq_le_parabolicInner, ShortTime.good_of_parabolic, ShortTime.rawJoint_limSec, ShortTime.jointlySmoothOn_limFam, ShortTime.jet_limFam, ShortTime.contDiffOn_limChart, ShortTime.ParabolicSpace.eq_zero_of_evalJet_eq_zero, ShortTime.ParabolicSpace.eq_ofSmooth_of_evalJet_eq}
\uses{lem:parabolic-basic, fact:sobolev}
Let $m\ge1$ and $i,j\ge0$ with $j\le m-1$ and $2p(m-j-1)>n/2+i$. Then
\[
  \sup_{M\times[0,T]}\abs{\partial_t^j\nabla^iu}\le C\norm{u}_{\Pp^m}\qquad\text{for smooth }u .
\]
Hence $\partial_t^j\nabla^iu$ is continuous for every $u\in\Pp^m$. If smooth $u_k$ form a Cauchy
sequence in $\Pp^m$ for every $m$, then all $\partial_t^j\nabla^iu_k$ converge uniformly on
$M\times[0,T]$, the limit $u$ is in $C^\infty(M\times[0,T];W)$, and
$\partial_t^j\nabla^iu=\lim_k\partial_t^j\nabla^iu_k$. In particular, if $u\in\Pp^m$ for all $m$, then
$u\in C^\infty(M\times[0,T];W)$.
\end{lemma}

\begin{proof}
Both $\partial_t^ju$ and $\partial_t^{j+1}u$ are bounded in $L^2(0,T;H^{2p(m-j-1)})$ by
$\norm u_{\Pp^m}$. Apply the inequality from the proof of
Lemma~\ref{lem:parabolic-basic}(v) with $X=H^{2p(m-j-1)}$, and then Fact~\ref{fact:sobolev}.

For the sequence, the estimate makes every $\partial_t^j\nabla^iu_k$ uniformly Cauchy. Fix a chart
and a local frame $B_1,\dots,B_n$ of vector fields with constant coordinates. In the chart and the
induced trivialization, the derivative of the reading of $u_k$ along the $\ell$-th coordinate is the
reading of $\nabla_{B_\ell}u_k$ minus a Christoffel term that is linear in $u_k$ with smooth
coefficients, and its time derivative is the reading of $\partial_tu_k$. These converge uniformly, so
the limit of the readings is $C^1$ with the limits as derivatives. The sequences
$\partial_tu_k$ and $\nabla_{B_\ell}u_k$ are again Cauchy in every $\Pp^m$, so induction gives
$C^\infty$. Evaluating at the centre of the chart identifies $\partial_tu$ and $\nabla_{B_\ell}u$ with
the limits of $\partial_tu_k$ and $\nabla_{B_\ell}u_k$. Since the $B_\ell$ span each tangent space, this
identifies $\nabla u$ with $\lim_k \nabla u_k$, and by induction all $\partial_t^j\nabla^iu$.

An element $v\in\Pp^m$ is determined by its continuous representative: if $v_k\to v$ with $v_k$ smooth
and $v_k\to0$ uniformly, then every $\partial_t^j\nabla^iv_k$ converges in $L^2$, and integrating by parts
against test sections vanishing at $t=0,T$ shows inductively that all these limits vanish, so $v=0$.

For the last claim, note that the inclusions $\Pp^{m+1}\hookrightarrow\Pp^m$ are compatible with the
embeddings into $L^2$. Apply the sequence statement to smooth approximations of $u$ in $\Pp^k$, $k\to\infty$.
\end{proof}

\begin{proposition}[algebra and composition]\label{prop:nemytskii}
\lean{ShortTime.ParabolicSpace.bilinear_le, ShortTime.ParabolicSpace.contDiff_comp, ShortTime.FaaDiBruno.exists_tame_bound_graph, ShortTime.FibreMap, ShortTime.FibreMap.exists_tame_pointwise, ShortTime.parabolicInner_le_of_jetSq_le, ShortTime.exists_jetSq_le_parabolicInner, ShortTime.SmoothExtension.exists_contDiffOn_extension, ShortTime.FibreMap.fibreDeriv, ShortTime.FibreMap.taylorMap, ShortTime.FibreMap.exists_block_bound, ShortTime.FibreMap.exists_derivFam_bound, ShortTime.FibreMap.exists_taylor_bound}
\uses{def:parabolic, lem:embedding}
Let $N\ge N_0:=\lfloor n/(2p)\rfloor+3$.
\begin{enumerate}[label=(\roman*),nosep]
\item $\sup\abs u\le C\norm u_{\Pp^N}$. Moreover, for every smooth bilinear bundle map
  $\beta\colon W_1\times W_2\to W_3$, possibly depending on $t$,
  \[
    \norm{\beta(u,v)}_{\Pp^N}\le C\norm u_{\Pp^N}\norm v_{\Pp^N}.
  \]
\item Let $E,E'$ be Riemannian bundles, $S\subseteq\R^d$ open, $\mathcal U\subseteq S\times[0,T]\times E$
  open, and $\Phi\colon\mathcal U\to E'$ smooth and fibre-preserving. Let
  $\mathcal V=\{(s,u)\in S\times\Pp^N(E):(s,t,u(x,t))\in\mathcal U\text{ for all }(x,t)\}$, which is
  open. Then the map
  \[
    \mathcal V\to\Pp^N(E'),\qquad(s,u)\mapsto\Phi_s(t,u)
  \]
  is $C^\infty$. Its derivative in $u$ is $h\mapsto D_u\Phi_s(t,u)h$.
\end{enumerate}
\end{proposition}

\begin{proof}
\emph{Step 1: a pointwise bound for low-order factors.} Let $2pj+i\le pN$. Then $j\le N-1$ and
\[
  2p(N-j-1)-i=2pN-2p-(2pj+i)\ge pN-2p>n/2
\]
by the choice of $N_0$. By Lemma~\ref{lem:embedding}, $\sup\abs{\partial_t^j\nabla^iu}\le C\norm u_{\Pp^N}$.

\emph{Step 2: the product estimate (i).} Expand $\partial_t^j\nabla^i\beta(u,v)$ with $2pj+i\le2pN$ by
the Leibniz rule. In each term the parabolic weights $2pj_1+i_1$ and $2pj_2+i_2$ of the two factors add
up to at most $2pN$, so one of them is at most $pN$. Estimate that factor in $L^\infty$ by Step~1 and
the other in $L^2$.

\emph{Step 3: composition.} Part (ii) is local, so we may work in charts and trivializations, with $s$
as additional variables that are constant in $(x,t)$. By the chain rule (Fa\`a di Bruno),
$\partial^\alpha_{t,x}[\Phi(t,x,u)]$ is a sum of terms
\[
  (\partial^\beta_{t,x}D^k_u\Phi)(t,x,u)\,[\partial^{\alpha_1}u,\dots,\partial^{\alpha_k}u],\qquad
  \beta+\alpha_1+\dots+\alpha_k=\alpha,\ \abs{\alpha_i}\ge1.
\]
The coefficients are bounded, since $u$ takes values in a compact subset of $\mathcal U$. At most one
$\alpha_i$ has parabolic weight larger than $pN$, so the terms are estimated as in Step~2. This shows
that $\Phi_*\colon u\mapsto\Phi(t,u)$ maps bounded subsets of $\mathcal V$ into bounded sets, first for
smooth $u$ and then, by continuity, for all $u$.

\emph{Step 4: differentiability.} Taylor's formula
\[
  \Phi(u+h)-\Phi(u)-D\Phi(u)h=\int_0^1(1-\tau)D^2\Phi(u+\tau h)[h,h]\,d\tau,
\]
combined with (i) and Step~3, gives a bound $C\norm h_{\Pp^N}^2$ on bounded sets. So $\Phi_*$ is
differentiable, and its derivative is multiplication by the $\Pp^N$-section $D\Phi(u)$. The map
$u\mapsto D\Phi(u)$ is again a composition operator, so by induction $\Phi_*$ is $C^\infty$. Joint
smoothness in $(s,u)$ follows in the same way.
\end{proof}

\begin{corollary}[the nonlinear operator]\label{cor:N}
\lean{ShortTime.EquationData.exists_nonlinOp, ShortTime.EquationData.exists_nonlinOp_aux, ShortTime.EquationData.exists_dom_tube, ShortTime.FibreMap.exists_nonlinOp_param, ShortTime.FibreMap.paramLift, ShortTime.FibreMap.exists_nonlinOp_param_ball}
\uses{def:data, def:linearization, def:parabolic}
Let $m\ge m_0=N_0+1$. The set
\[
  \Pp^m_\Dom=\{u\in\Pp^m(V):\ (t,j^{2p}u(x,t))\in\Dom\text{ for all }(x,t)\}
\]
is open. The map $u\mapsto F[u]$ is $C^\infty$ from $\Pp^m_\Dom$ to $\Pp^{m-1}(V)$, with derivative
\[
  v\mapsto\sum_kD_kF(t,j^{2p}u)(\nabla^kv).
\]
The same holds for smooth families $F_s$ depending on a parameter $s$, jointly in $(s,u)$.

In the formalization the statement is local: near a smooth $w$ with jets in $\Dom$, the map is
defined on a ball of $\Pp^m$ around $w$, agrees with $u\mapsto F[u]$ on smooth $u$, is $C^\infty$ at
$w$ (for each $n$ it is $C^n$ on a ball whose radius depends on $n$), and its derivative at $w$ is
$v\mapsto\sum_kD_kF(t,j^{2p}w)(\nabla^kv)$. This is what the implicit function
theorem needs.
\end{corollary}

\begin{proof}
\uses{prop:nemytskii, lem:parabolic-basic}
Compose the bounded linear map $j^{2p}\colon\Pp^m\to\Pp^{m-1}(\Jet^{2p}V)$ of
Lemma~\ref{lem:parabolic-basic}(i) with Proposition~\ref{prop:nemytskii}, using $N=m-1\ge N_0$. In the formalization the derivatives $g_k(u)=D^k_\zeta F(t,j^{2p}u)[j^{2p}h_1,\dots,j^{2p}h_k]$
satisfy, for smooth $u$ near $w$ and smooth $h_i$, the bounds
$\norm{g_k(u)(h)}\le C\prod\norm{h_i}$ and
$\norm{g_k(u+h)-g_k(u)-g_{k+1}(u)(h,\cdot)}\le C\norm h^2$ (Taylor's formula with integral
remainder at fibre level, and Proposition~\ref{prop:nemytskii}); a map with such Taylor bounds on a
dense subspace extends to a $C^n$ map on the completion.
\end{proof}

\begin{lemma}[spatial composition]\label{lem:spatial-nemytskii}
\lean{ShortTime.EquationData.contDiff_spatialOp}
\uses{def:data, def:sobolev, fact:sobolev}
Let $k>n/2$. The map $a\mapsto F(t,j^{2p}a)$ is $C^\infty$ from $H^{k+2p}(M;V)$ to $H^k(M;V)$, jointly
in $t$. It is defined on the open set of those $a$ whose jets lie in the slice $\Dom_t$.
\end{lemma}

\begin{proof}
The proof of Proposition~\ref{prop:nemytskii}, without time derivatives.
\end{proof}

\subsection{Spectral facts}

\begin{fact}[eigenbasis {\cite{Aubin}}]\label{fact:eigen}
\source{DG: Analysis/Spectral/Tensor (TensorEigenIdx; tensor bundles)}
There is an $L^2$-orthonormal basis $(\varphi_i)_{i\ge1}$ of $L^2(M;W)$ consisting of smooth sections
with $\Lap\varphi_i=\mu_i\varphi_i$, where $0\le\mu_1\le\mu_2\le\cdots\to\infty$.
\end{fact}

\begin{fact}[elliptic estimate]\label{fact:elliptic}
\uses{def:sobolev}
\source{DG: Analysis/Spectral/Tensor/SobolevScale/Jet/Bounds/IteratedCovariantDerivative (hsJet\_le, hs\_le\_jet; covariant tensors of every rank)}
For every $k\in\N$ there is $C$ such that
\[
  \norm{u}_{H^{2k}}\le C\bigl(\norm{(\Lap)^ku}+\norm u\bigr)\qquad\text{for smooth }u.
\]
This follows from Bochner-type integrations by parts together with Lemma~\ref{lem:ricci}, by
induction on $k$.
\end{fact}

Let $\Pi_N$ denote the $L^2$-orthogonal projection onto the span of $\varphi_1,\dots,\varphi_N$. Since
$\Lap$ is symmetric on smooth sections, $\Pi_N(\Lap)^k=(\Lap)^k\Pi_N$ on smooth sections.

% ===========================================================================
\section{Linear theory}\label{sec:linear}
% ===========================================================================

\subsection{The class of operators and G\aa rding's inequality}

\begin{definition}\label{def:class}
\lean{ShortTime.LinearOp, ShortTime.LinearOp.apply, ShortTime.LinearOp.Solves, ShortTime.LinearOp.IsParabolic}
\uses{def:legendre}
Let $W$ be a Riemannian bundle, $T\in(0,T_0]$, $\lambda>0$, $K\in\N$ and $\Lambda\ge0$. The class
$\Lcl(W;T,\lambda,K,\Lambda)$ consists of the operators
\[
  Lu=P(\nabla^{2p}u)+\sum_{k<2p}R_k(\nabla^ku)
\]
with the following properties:
\begin{itemize}[nosep]
\item the coefficients $P$ and $R_k$ are smooth on $M\times[0,T]$;
\item $\kappa(P(x,t))\ge\lambda$ for all $(x,t)$;
\item the sum of $\sup\abs{\partial_t^a\nabla^bP}$ and $\sup\abs{\partial_t^a\nabla^bR_k}$ over
  $a+b\le K$ and all $k$ is at most $\Lambda$.
\end{itemize}
\end{definition}

\begin{definition}[flip of a coefficient]\label{def:flip}
\lean{ShortTime.flipField, ShortTime.flipRemainder}
Let $E,F$ be Riemannian bundles and $A\in\Gamma(\Hom(T^*M\otimes E,F))$, the first slot being the
derivative slot. Its \emph{flip} is $\check A\in\Gamma(\Hom(E,T^*M\otimes F))$,
\[
  \check A(w)(v,\cdot)=A(v^\flat\otimes w)(\cdot),\qquad v^\flat=g(v,\cdot).
\]
Its \emph{remainder} is $R_A=\operatorname{tr}_{12}\circ\nabla\check A\in\Gamma(\Hom(E,F))$, where
$\operatorname{tr}_{12}$ contracts the two leading slots with $g^{-1}$.
\end{definition}

\begin{lemma}[one integration by parts]\label{lem:ibp-one}
\lean{ShortTime.inner_apply_covGrad, ShortTime.cometricDoubleTraceFib_comp_slotExtend_flipFib, ShortTime.covDivergence_eq_trace_covGrad}
\uses{def:flip}
For $T\in C^\infty(M;F)$ and $W\in C^\infty(M;E)$,
\[
  \ip{T}{A(\nabla W)}=-\ip{\nabla T}{\check A\,W}-\ip{T}{R_A W}.
\]
\end{lemma}

\begin{proof}
In an orthonormal frame, $\operatorname{tr}_{12}\bigl((\mathrm{id}_{T^*M}\otimes\check A)D\bigr)
=\sum_iA(e_i^\flat\otimes D(e_i,\cdot))=A(D)$ for every $D\in T^*M\otimes T^*M\otimes E$, because
$\sum_ie_i^\flat\otimes D(e_i,\cdot)=D$. By the Leibniz rule
$\nabla(\check AW)=(\nabla\check A)W+(\mathrm{id}\otimes\check A)(\nabla W)$, so
$A(\nabla W)=\operatorname{tr}_{12}\nabla(\check AW)-R_AW$. Finally
$\operatorname{tr}_{12}\nabla V=\operatorname{div}V$ and $\ip{T}{\operatorname{div}V}=-\ip{\nabla T}{V}$.
\end{proof}

\begin{lemma}[integration by parts in the principal term]\label{lem:ibp}
\lean{ShortTime.principal_ibp, ShortTime.principal_iter, ShortTime.evalFirst_flipIter, ShortTime.legendreForm_eq_fiberInner_flipIter}
\uses{def:legendre, def:sobolev, lem:ibp-one}
Let $P$ be a smooth section of $\Hom(\otimes^{2p}T^*M\otimes W,W)$ and $\psi\in C^\infty(M;W)$. Then
\[
  \ip{\psi}{P(\nabla^{2p}\psi)}=(-1)^p\ip{\nabla^p\psi}{\hatP\,\nabla^p\psi}+\mathcal E(\psi),
  \qquad
  \abs{\mathcal E(\psi)}\le C\,N_{p-1}(\psi)\,N_p(\psi),
\]
where $C$ depends on $\sup_{a\le p}\abs{\nabla^aP}$ and on the curvature. If $P$ depends smoothly on
$t\in[0,T]$, then $C$ can be chosen uniformly in $t$.
\end{lemma}

\begin{proof}
\uses{lem:ricci-l2}
Apply Lemma~\ref{lem:ibp-one} $p$ times, each time moving the outermost derivative of
$\nabla^{2p-j}\psi=\nabla(\nabla^{2p-j-1}\psi)$ onto $\nabla^j\psi$. With $P^{(0)}=P$ and
$P^{(j+1)}=\check{P^{(j)}}$ this gives
\[
  \ip{\psi}{P(\nabla^{2p}\psi)}=(-1)^p\ip{\nabla^p\psi}{P^{(p)}\nabla^p\psi}
  -\sum_{j<p}(-1)^j\ip{\nabla^j\psi}{R_{P^{(j)}}\nabla^{2p-j-1}\psi}.
\]
Unwinding the definition, $P^{(p)}(\xi)(v_1,\dots,v_p,\cdot)=P(v_p^\flat\otimes\dots\otimes v_1^\flat
\otimes\xi)$. Expanding in an orthonormal frame over the first $p$ slots and reindexing $I$ by its
reversal gives $\ip{\rho\xi}{P^{(p)}\eta}=\ip{\xi}{\hatP\eta}$ pointwise, where $\rho$ reverses the
order of the first $p$ slots. Hence
\[
  \ip{\nabla^p\psi}{P^{(p)}\nabla^p\psi}=\ip{\nabla^p\psi}{\hatP\,\nabla^p\psi}
  +\ip{\nabla^p\psi-\rho\nabla^p\psi}{P^{(p)}\nabla^p\psi},
\]
and the last term is at most $C\,N_{p-2}(\psi)\norm{\nabla^p\psi}$ by Lemma~\ref{lem:ricci-l2}. (For
$p=1$, $\rho$ is the identity.)

\emph{The error terms.} Each $R_{P^{(j)}}$ has order zero and involves one derivative of $P$. In the
term $\ip{\nabla^j\psi}{R\nabla^{2p-j-1}\psi}$ move $p-j-1$ further derivatives from the right to
the left, again by Lemma~\ref{lem:ibp-one}. This produces pairings of $\nabla^a\psi$ with
$\nabla^b\psi$, $a\le p$, $b\le p-1$, whose coefficients involve at most $p$ derivatives of $P$.
Each is at most $C\,N_{p-1}(\psi)N_p(\psi)$.
\end{proof}

\begin{lemma}[G\aa rding's inequality]\label{lem:garding}
\lean{ShortTime.LinearOp.garding, ShortTime.LinearOp.garding_top, ShortTime.LinearOp.coercive}
\uses{def:class, def:sobolev}
Let $L\in\Lcl(W;T,\lambda,K,\Lambda)$ with $K\ge p$. There is $C=C(n,p,\lambda,\Lambda,M,W)$ such that
for all $t\in[0,T]$ and all $\psi\in C^\infty(M;W)$
\[
  -\ip{\psi}{L(t)\psi}\ge\frac\lambda4\norm\psi_{H^p}^2-C\norm\psi^2 .
\]
\end{lemma}

\begin{proof}
\uses{lem:ibp, lem:interp}
Write $N_k=N_k(\psi)$ as in Lemma~\ref{lem:interp}.

\emph{The principal term.} By Lemma~\ref{lem:ibp} and the Legendre condition
$(-1)^{p+1}\ip{\xi}{\hatP\xi}\ge\lambda\abs\xi^2$, applied pointwise to $\xi=\nabla^p\psi$,
\[
  -\ip{\psi}{P(\nabla^{2p}\psi)}\ge\lambda\norm{\nabla^p\psi}^2-C N_{p-1}N_p.
\]

\emph{Lower-order terms of order $k\le p$.} Here
$\abs{\ip{\psi}{R_k\nabla^k\psi}}\le\Lambda N_{p-1}N_p$.

\emph{Lower-order terms of order $p<k<2p$.} Write $\nabla^k\psi=\nabla^{k-p}(\nabla^p\psi)$ and
integrate by parts $k-p$ times. This produces terms $\ip{\nabla^b R_k*\nabla^a\psi}{\nabla^p\psi}$
with $a+b=k-p\le p-1$, which are bounded by $C N_{p-1}N_p$.

\emph{Conclusion.} So $\ip{\psi}{L\psi}\le-\lambda\norm{\nabla^p\psi}^2+C N_{p-1}N_p$. Use
$N_p\le N_{p-1}+\norm{\nabla^p\psi}$ and $N_{p-1}\le\delta\norm{\nabla^p\psi}+C_\delta\norm\psi$ from
Lemma~\ref{lem:interp} with $\delta$ small, and absorb:
$\ip{\psi}{L\psi}+\frac\lambda2\norm{\nabla^p\psi}^2\le C\norm\psi^2$. Finally
$\norm\psi_{H^p}^2\le2\norm{\nabla^p\psi}^2+C\norm\psi^2$.
\end{proof}

\subsection{Energy estimates}

For $r\in\N$ let $W_r=\otimes^rT^*M\otimes V$, so that $W_0=V$ and $W_{r+1}=T^*M\otimes W_r$.

\begin{lemma}[differentiating an operator]\label{lem:commute}
\lean{ShortTime.LinearOp.exists_covGrad_comm, ShortTime.LinearOp.isParabolic_comm, ShortTime.LinearOp.legendreForm_comm, ShortTime.isDiffOp_cycle_defect, ShortTime.IsDiffOp.split}
\uses{def:class}
Let $L\in\Lcl(W_r;T,\lambda,K,\Lambda)$ with $K\ge1$. There are an operator
$\tilde L\in\Lcl(W_{r+1};T,\lambda,K-1,\tilde\Lambda)$ and a smooth
$S_0\in\Gamma(\Hom(W_r,W_{r+1}))$, depending on $t$, with
\[
  \nabla(Lu)=\tilde L(\nabla u)+S_0u\qquad\text{for all }u.
\]
The constant $\tilde\Lambda$ and the $C^{K-1}$-norm of $S_0$ are bounded in terms of $\Lambda$, $r$
and the geometry.
\end{lemma}

\begin{proof}
\uses{lem:ricci, lem:legendre-basic}
\emph{The new principal coefficient.} Let $P^{(1)}\colon\otimes^{2p}T^*M\otimes W_{r+1}\to W_{r+1}$
act by $P$ on the $2p$ form slots and on the $W_r$ factor, leaving the extra $T^*M$ factor untouched.
The first slot of $\nabla u\in\Gamma(T^*M\otimes W_r)$ is the derivative slot, and
$\nabla^{2p}(\nabla u)=\nabla^{2p+1}u$.

\emph{Comparing the two expressions.} By the Leibniz rule,
\[
  \nabla\bigl(P(\nabla^{2p}u)\bigr)=(\nabla P)(\nabla^{2p}u)+P_\bullet(\nabla^{2p+1}u),
\]
where $P_\bullet$ applies $P$ to slots $2,\dots,2p+1$. Moving the first slot to position $2p+1$ is a cyclic
permutation $\gamma$, so by Lemma~\ref{lem:ricci}
\[
  P_\bullet(\nabla^{2p+1}u)=P^{(1)}\bigl(\nabla^{2p}(\nabla u)\bigr)+P^{(1)}(Q_\gamma u),\qquad
  Q_\gamma\in\Diff^{2p-1}.
\]
Everything else, including $\nabla(R_k\nabla^ku)$, lies in $\Diff^{2p}(W_r,W_{r+1})$.

\emph{Rewriting in terms of $\nabla u$.} Every $A\in\Diff^{2p}(W_r,W_{r+1})$ can be written as
$Au=A_0u+A'(\nabla u)$ with $A'\in\Diff^{2p-1}(W_{r+1},W_{r+1})$, since
$\nabla^ku=\nabla^{k-1}(\nabla u)$. Set $\tilde L=P^{(1)}\nabla^{2p}+A'$ and $S_0=A_0$.

\emph{The Legendre condition.} Decompose $\xi\in\otimes^pT^*M\otimes W_{r+1}$ as
$\xi=\sum_c\xi^{(c)}\otimes e^c$ along the extra factor. Then
$\ip{\xi}{\widehat{P^{(1)}}\xi}=\sum_c\ip{\xi^{(c)}}{\hatP\xi^{(c)}}$, so
$\kappa(P^{(1)})\ge\kappa(P)\ge\lambda$.
\end{proof}

\begin{proposition}[energy estimates]\label{prop:energy}
\lean{ShortTime.energy_estimate, ShortTime.LinearOp.energy_estimate_zero_of_garding, ShortTime.LinearOp.garding_lerp, ShortTime.hasDerivWithinAt_covGrad, ShortTime.energy_ode}
\uses{def:class, def:sobolev}
Let $K_k=p+k$. For every $k\in\N$ and $r\in\N$, every $L\in\Lcl(W_r;T,\lambda,K_k,\Lambda)$, and every
$u\in C^\infty(M\times[0,T];W_r)$ with $u(\cdot,0)=0$, set $f=\partial_tu-Lu$. Then
\begin{equation}\label{eq:Ek}
  \max_{t\in[0,T]}\norm{u(t)}_{H^k}^2+\int_0^T\norm{u}_{H^{k+p}}^2dt\le C\int_0^T\norm{f}_{H^{k-p}}^2dt,
\end{equation}
with $C=C(k,r,n,p,\lambda,\Lambda,T_0,M,V)$.
\end{proposition}

\begin{proof}
\uses{lem:garding, lem:commute, lem:sobolev-basic, fact:gronwall}
We argue by induction on $k$, for all $r$ simultaneously.

\emph{Base case $k=0$.} By Lemma~\ref{lem:garding} and
Lemma~\ref{lem:sobolev-basic}(ii),
\begin{align*}
\frac{d}{dt}\norm u^2&=2\ip u{Lu}+2\ip uf
\le-\frac\lambda2\norm u_{H^p}^2+2C\norm u^2+2\norm u_{H^p}\norm f_{H^{-p}}\\
&\le-\frac\lambda4\norm u_{H^p}^2+2C\norm u^2+\frac4\lambda\norm f_{H^{-p}}^2 .
\end{align*}
Since $u(0)=0$, Gronwall's inequality gives
$\max_t\norm{u(t)}^2\le\frac4\lambda e^{2CT_0}\int_0^T\norm f_{H^{-p}}^2$. Integrating the differential
inequality over $[0,T]$ then bounds $\int_0^T\norm u_{H^p}^2$.

\emph{Induction step from $k$ to $k+1$.} Let $L\in\Lcl(W_r;T,\lambda,K_{k+1},\Lambda)$ and $w=\nabla u$,
so that $w(\cdot,0)=0$. By Lemma~\ref{lem:commute},
\[
  \partial_tw=\tilde Lw+S_0u+\nabla f,\qquad\tilde L\in\Lcl(W_{r+1};T,\lambda,K_k,\tilde\Lambda).
\]
Applying the induction hypothesis to $w$, followed by Lemma~\ref{lem:sobolev-basic}(iii)--(iv), gives
\[
  \max_t\norm{w}_{H^k}^2+\int_0^T\norm w_{H^{k+p}}^2
  \le C\int_0^T\Bigl(\norm u^2_{H^{\max(k-p,0)}}+\norm f_{H^{k+1-p}}^2\Bigr).
\]
Next,
\[
  \int_0^T\norm u^2_{H^{\max(k-p,0)}}\le T_0\max_t\norm u_{H^k}^2
  \le C\int_0^T\norm f_{H^{k-p}}^2\le C\int_0^T\norm f^2_{H^{k+1-p}},
\]
using the induction hypothesis for $u$ (valid because $K_{k+1}\ge K_k$) and the monotonicity in
Lemma~\ref{lem:sobolev-basic}(ii). Finally
$\norm u_{H^{k+1}}^2=\norm u^2+\norm w_{H^k}^2$ and
$\norm u_{H^{k+1+p}}^2=\norm u^2+\norm w_{H^{k+p}}^2$ by Lemma~\ref{lem:sobolev-basic}(i). Combined
with the case $k=0$, this gives \eqref{eq:Ek} for $k+1$.
\end{proof}

\begin{remark}[uniformity in the formalization]
In the formalization the constant is not tracked through the class $\Lcl$. Instead,
Proposition~\ref{prop:energy} is proved, by the same induction, for all operators
$L_\theta=(1-\theta)L_0+\theta L_1$, $\theta\in[0,1]$, on the segment between two operators $L_0,L_1$ that
are parabolic with the same constant $\lambda$, with one constant $C$ for the whole segment. This is
the uniformity used in the proof of Theorem~\ref{thm:linear}. The base case needs G\aa rding's inequality
uniformly in $\theta$; it holds because $\ip u{L_\theta u}$ is affine in $\theta$. In the induction step,
Lemma~\ref{lem:commute} is applied to $L_0$ and $L_1$ separately, and $\tilde L_\theta=(1-\theta)\tilde L_0+\theta\tilde L_1$.
The identity $\partial_t\nabla u=\nabla\partial_tu$ for smooth families is proved weakly, by pairing with test
fields and integrating by parts. Proposition~\ref{prop:apriori} is proved in the same form,
uniformly along such segments.
\end{remark}

\begin{remark}
Negative-order norms enter only through the forcing term. They let the induction start at $k=0$ and
reach the ``maximal regularity'' level $k=p$ by differentiation alone. No separate argument is needed
for the second-order-in-time energy.
\end{remark}

\begin{proposition}[a priori estimate]\label{prop:apriori}
\lean{ShortTime.apriori_estimate, ShortTime.parabolicInner_self_eq, ShortTime.LinearOp.exists_sobolevNormSq_timeDerivIter_apply_le, ShortTime.timeDerivIter_iteratedCovGrad}
\uses{def:class, def:parabolic}
Let $m\ge1$ and $L\in\Lcl(V;T,\lambda,K,\Lambda)$ with $K\ge(2p+1)m$. Then
\[
  \norm u_{\Pp^m}\le C\norm{(\partial_t-L)u}_{\Pp^{m-1}}\qquad\text{for all }u\in\Pp^m_0(V),
\]
with $C=C(m,n,p,\lambda,\Lambda,T_0,M,V)$.
\end{proposition}

\begin{proof}
\uses{prop:energy, lem:parabolic-basic, lem:sobolev-basic}
By density and Lemma~\ref{lem:parabolic-basic}(iii) it suffices to treat smooth $u$ with $u(0)=0$. Let
$f=\partial_tu-Lu$.

\emph{Top spatial regularity.} Proposition~\ref{prop:energy} with $k=2pm-p$ gives
\[
  \int_0^T\norm u^2_{H^{2pm}}\le C\int_0^T\norm f^2_{H^{2p(m-1)}}\le C\norm f_{\Pp^{m-1}}^2 .
\]

\emph{Time derivatives.} For $1\le j\le m$ we have
\[
  \partial_t^ju=\sum_{a=0}^{j-1}\binom{j-1}{a}(\partial_t^aL)(\partial_t^{j-1-a}u)+\partial_t^{j-1}f,
\]
where $\partial_t^aL$ has the $\partial_t^a$-differentiated coefficients. Set $b=j-1-a\le j-1$. By
Lemma~\ref{lem:sobolev-basic}(iv),
\[
  \norm{(\partial_t^aL)\partial_t^bu}_{H^{2p(m-j)}}\le C\norm{\partial_t^bu}_{H^{2p(m-j)+2p}}
  \le C\norm{\partial_t^bu}_{H^{2p(m-b)}} .
\]
Induction on $j$ therefore bounds $\int_0^T\norm{\partial_t^ju}_{H^{2p(m-j)}}^2$ by
$C\norm f_{\Pp^{m-1}}^2$.
\end{proof}

\subsection{The model operator}

\begin{proposition}\label{prop:model}
\lean{ShortTime.LinearOp.exists_model, ShortTime.exists_polyharmonic, ShortTime.LinearOp.exists_solution_proj, ShortTime.LinearOp.approx_of_eigen, ShortTime.tendsto_integral_proj_sub, ShortTime.LinearOp.surjective_heatExt_of_approx}
\uses{def:class, def:parabolic}
The operator $L_0=-(\Lap)^p$ on $V$ belongs to $\Lcl(V;T,1,K,\Lambda_0(K))$ for every $K$. For every
$m\ge1$, the map $\partial_t-L_0\colon\Pp^m_0(V)\to\Pp^{m-1}(V)$ is surjective.
\end{proposition}

\begin{proof}
\uses{lem:legendre-examples, fact:eigen, fact:elliptic, prop:apriori}
\emph{Membership in the class.} This is Lemma~\ref{lem:legendre-examples}(b).

\emph{Solving mode by mode.} Let $f\in C^\infty(M\times[0,T];V)$ and $f_i(t)=\ip{f(\cdot,t)}{\varphi_i}$,
with $\varphi_i$ as in Fact~\ref{fact:eigen}. Set
\[
  c_i(t)=\int_0^te^{-\mu_i^p(t-\tau)}f_i(\tau)\,d\tau .
\]
Then $c_i$ is smooth, $c_i(0)=0$ and $c_i'=-\mu_i^pc_i+f_i$. Hence $u_N=\sum_{i\le N}c_i\varphi_i$ is
smooth, $u_N(0)=0$, and $(\partial_t-L_0)u_N=\Pi_Nf$.

\emph{Density.} We claim that $\Pi_Nf\to f$ in $\Pp^{m-1}$. Indeed $\partial_t^j\Pi_Nf=\Pi_N\partial_t^jf$, and for a
smooth section $h$ of $V$, Fact~\ref{fact:elliptic} gives
\[
  \norm{h-\Pi_Nh}_{H^{2k}}^2\le C\sum_{i>N}(1+\mu_i)^{2k}\ip{h}{\varphi_i}^2 .
\]
The right-hand side tends to $0$ and is bounded by $C\norm h_{H^{2k}}^2$. Apply this with
$h=\partial_t^jf(\cdot,t)$ and use dominated convergence in $t$. Since $(\partial_t-L_0)u_N=\Pi_Nf$
with $u_N$ smooth and $u_N(0)=0$, and smooth sections are dense in $\Pp^{m-1}$, the range of
$\partial_t-L_0\colon\Pp^m_0\to\Pp^{m-1}$ is dense.

\emph{Surjectivity.} By Proposition~\ref{prop:apriori}, $\norm u_{\Pp^m}\le C\norm{(\partial_t-L_0)u}_{\Pp^{m-1}}$
on $\Pp^m_0$, which is complete, so the range is closed. Hence it is all of $\Pp^{m-1}$.
\end{proof}

\subsection{The method of continuity and the linear theorem}

\begin{lemma}[method of continuity {\cite[Thm.~5.2]{GT}}]\label{lem:continuity}
\lean{ShortTime.surjective_of_continuity}
Let $X$ be a Banach space, $Y$ a normed space and $T_0,T_1\in\mathcal B(X,Y)$. Put
$T_s=(1-s)T_0+sT_1$, and suppose that $\norm x\le C\norm{T_sx}$ for all $s\in[0,1]$ and $x\in X$. If
$T_0$ is surjective, then so is $T_1$.
\end{lemma}

\begin{proof}
Suppose $T_\sigma$ is surjective. By the hypothesis it is then invertible, with
$\norm{T_\sigma^{-1}}\le C$. For $\abs{s-\sigma}<\bigl(C(\norm{T_0}+\norm{T_1})\bigr)^{-1}$, the
equation $T_sx=y$ is equivalent to the fixed-point problem
\[
  x=T_\sigma^{-1}y-(s-\sigma)T_\sigma^{-1}(T_1-T_0)x,
\]
whose right-hand side is a contraction. Hence $T_s$ is surjective. Steps of this fixed length cover
$[0,1]$.
\end{proof}

\begin{theorem}[linear parabolic theory]\label{thm:linear}
\lean{ShortTime.LinearOp.exists_smooth, ShortTime.LinearOp.unique_of_isParabolic, ShortTime.LinearOp.isomorphism, ShortTime.LinearOp.isomorphism_of_base, ShortTime.LinearOp.heatExt, ShortTime.LinearOp.norm_le_heatExt_lerp, ShortTime.LinearOp.exists_approx, ShortTime.LinearOp.exists_smooth_zero}
\uses{def:class, def:parabolic}
Let $m\ge1$ and $L\in\Lcl(V;T,\lambda,K,\Lambda)$ with $K\ge(2p+1)m$.
\begin{enumerate}[label=(\roman*),nosep]
\item The map $\partial_t-L\colon\Pp^m_0(V)\to\Pp^{m-1}(V)$ is a linear homeomorphism.
\item Suppose $L$ has smooth coefficients and let $f\in C^\infty(M\times[0,T];V)$ and
  $u_0\in C^\infty(M;V)$. Then there is exactly one $u\in C^\infty(M\times[0,T];V)$ with
  $\partial_tu=Lu+f$ and $u(0)=u_0$.
\end{enumerate}
\end{theorem}

\begin{proof}
\uses{prop:model, prop:apriori, lem:continuity, lem:legendre-basic, lem:embedding, lem:garding}
(i) Let $L_s=(1-s)L_0+sL$. By Lemma~\ref{lem:legendre-basic}(i),
$L_s\in\Lcl(V;T,\min(1,\lambda),K,\max(\Lambda_0,\Lambda))$ for all $s\in[0,1]$, so
Proposition~\ref{prop:apriori} holds uniformly in $s$. Lemma~\ref{lem:continuity} and
Proposition~\ref{prop:model} give surjectivity. Injectivity and continuity of the inverse follow from
Proposition~\ref{prop:apriori}.

(ii) Put $v=u-u_0$, with $u_0$ constant in time. We must find a smooth $v$ with $v(0)=0$ and
$(\partial_t-L)v=h$, where $h=f+Lu_0$ is smooth.
\begin{itemize}[nosep]
\item For each $k$, (i) at level $m=k+1$ gives $V\in\Pp^{k+1}_0$ with $(\partial_t-L)V=h$. Approximating
  $V$ by smooth families vanishing at $t=0$ gives a smooth $v_k$ with $v_k(0)=0$ and
  $\norm{(\partial_t-L)v_k-h}^2_{\Pp^k}<1/(k+1)$.
\item By Proposition~\ref{prop:apriori} at level $q+1$, for $k,k'\ge q$,
  $\norm{v_k-v_{k'}}^2_{\Pp^{q+1}}\le C\norm{(\partial_t-L)(v_k-v_{k'})}^2_{\Pp^q}
  \le 2C\bigl(\tfrac1{k+1}+\tfrac1{k'+1}\bigr)$. So $(v_k)$ is Cauchy in every $\Pp^q$.
\item By Lemma~\ref{lem:embedding}, $v_k$ converges with all derivatives $\partial_t^j\nabla^i$,
  uniformly, to a smooth $v$ with $v(0)=0$. Likewise $(\partial_t-L)v_k\to h$ uniformly.
\item Passing to the limit in $\partial_tv_k=Lv_k+(\partial_t-L)v_k$ gives $\partial_tv=Lv+h$.
\end{itemize}
For uniqueness, the difference $w$ of two smooth solutions is smooth, vanishes at $t=0$ and solves
$\partial_tw=Lw$. By Lemma~\ref{lem:garding},
$\frac{d}{dt}\norm{w}^2=2\ip{w}{Lw}\le2C\norm w^2$, so $w=0$ by Gr\"onwall's inequality. (Alternatively,
$w\in\Pp^m_0$ is zero by (i).)
\end{proof}

% ===========================================================================
\section{The nonlinear theory}\label{sec:nonlinear}
% ===========================================================================

Throughout this section $F$, $\Dom$ and $\Par$ are as in Section~\ref{sec:setting}, and $m\ge m_0$ is
fixed.

\subsection{Approximate solutions}

\begin{lemma}[formal solution]\label{lem:formal}
\lean{ShortTime.EquationData.formalCoeff, ShortTime.EquationData.formalCoeff_spec, ShortTime.EquationData.formalCoeff_spec', ShortTime.EquationData.formalPoly, ShortTime.EquationData.timeDerivIter_formalPoly, ShortTime.EquationData.exists_formalPoly_dom, ShortTime.EquationData.timeDerivIter_nonlinFamily_congr}
\uses{def:data}
Let $u_0\in C^\infty(M;V)$ with $(0,j^{2p}u_0(x))\in\Dom$ for all $x$. There are sections
$a_0=u_0,a_1,a_2,\ldots\in C^\infty(M;V)$ with the following property. Let $N\ge1$, $T'>0$, and let
$w$ be a smooth section on $M\times[0,T']$ with jets in $\Dom$ and $\partial_t^\ell w(0)=a_\ell$ for
$\ell\le N$. Then
\[
  \partial_t^\ell\big|_{t=0}\bigl(\partial_tw-F[w]\bigr)=0\qquad\text{for }0\le\ell\le N-1 .
\]
\end{lemma}

\begin{proof}
Since $\partial_t$ commutes with $\nabla$, we have $\partial_t^i j^{2p}w=j^{2p}\partial_t^iw$. By the
chain rule there are smooth fibre maps $\Phi_\ell$, determined by $F$ and its derivatives at $t=0$,
such that
\[
  \partial_t^\ell\big|_{t=0}F[w]=\Phi_\ell\bigl(j^{2p}a_0,\dots,j^{2p}a_\ell\bigr)
\]
for every such $w$: by the chain rule (Fa\`a di Bruno), $\partial_t^\ell|_{t=0}F[w]$ depends only on
$\partial_t^iw(0)$, $i\le\ell$. Define recursively $a_{\ell+1}=\partial_t^\ell|_{t=0}F[w_\ell]$ with
$w_\ell=\sum_{i\le\ell}a_it^i/i!$, taken on a short interval on which the jets of $w_\ell$ lie in $\Dom$;
such an interval exists since $\Dom$ is open and $M$ is compact, and one-sided derivatives at $t=0$ do
not depend on the interval. Then
$\partial_t^\ell\partial_tw(0)=a_{\ell+1}=\partial_t^\ell F[w](0)$ whenever $\ell+1\le N$.
\end{proof}

\begin{lemma}[cutoff]\label{lem:cutoff}
\lean{ShortTime.ParabolicSpace.tendsto_cutoff, ShortTime.cutoffFamily, ShortTime.jointlySmoothOn_cutoffFamily, ShortTime.cutoffFamily_eq_zero, ShortTime.tendsto_parabolicInner_cutoff}
\uses{def:parabolic}
Let $f\in C^\infty(M\times[0,T];V)$ with $\partial_t^\ell f(\cdot,0)=0$ for $\ell\le N-1$. Let
$\chi\in C^\infty(\R)$ with $\chi=0$ on $(-\infty,1]$ and $\chi=1$ on $[2,\infty)$, and set
$f_\varepsilon(x,t)=\chi(t/\varepsilon)f(x,t)$. Then $f_\varepsilon$ is smooth, vanishes for
$t\le\varepsilon$, and $f_\varepsilon\to f$ in $\Pp^{m'}$ as $\varepsilon\to0$, for every $m'\le N$.
\end{lemma}

\begin{proof}
\emph{A pointwise bound.} By Taylor's theorem, $\abs{\partial_t^l\nabla^if(x,t)}\le C_it^{N-l}$ for
$l\le N$.

\emph{Expanding the difference.} We have
\[
  \partial_t^j\nabla^i(f-f_\varepsilon)
  =\sum_{q\le j}\binom jq\,\partial_t^q\bigl(1-\chi(t/\varepsilon)\bigr)\,\partial_t^{j-q}\nabla^if,
\]
which is supported in $t\le2\varepsilon$.
\begin{itemize}[nosep]
\item The terms with $q\ge1$ are bounded by
  $C\varepsilon^{-q}(2\varepsilon)^{N-j+q}=C\varepsilon^{N-j}$, which is at most $C$ for $j\le N$.
\item The term with $q=0$ is bounded.
\end{itemize}
Since the support has measure $O(\varepsilon)$, all $L^2$ norms are $O(\varepsilon^{1/2})$.
\end{proof}

\subsection{Existence}\label{sec:existence}

Let $u_0$ satisfy the hypothesis of Theorem~\ref{thm:A}. Take the sections $a_\ell$ of
Lemma~\ref{lem:formal} and set
\[
  w(x,t)=\sum_{\ell\le m}a_\ell(x)\frac{t^\ell}{\ell!}.
\]
Then $j^{2p}w(\cdot,0)=j^{2p}u_0$. By Lemma~\ref{lem:legendre-basic}(iii) and continuity there are
$T_1\in(0,T_0]$ and $\lambda_1>0$ such that $w$ has jets in $\Par$ on $[0,T_1]$ and
$\kappa(D_{2p}F(t,j^{2p}w))\ge\lambda_1$ there. Let $\delta_1>0$ be such that the $\delta_1$-tube
(in the sense of Lemma~\ref{lem:legendre-basic}(iii)) around the compact set
$K_w=\{(t,j^{2p}w(x,t))\}$ is contained in $\Par$.

From now on $T=T_1$ in all function spaces. Set $f=\partial_tw-F[w]$; it is smooth, and
$\partial_t^\ell f(0)=0$ for $\ell\le m-1$ by Lemma~\ref{lem:formal}. Define
\[
  G(v)=\partial_t(w+v)-F[w+v]\qquad\text{on}\qquad\mathcal U=\{v\in\Pp^m_0:\ w+v\in\Pp^m_\Dom\}.
\]
By Corollary~\ref{cor:N}, $G$ is $C^\infty$ from $\mathcal U$ to $\Pp^{m-1}$, with $G(0)=f$ and
$DG(0)=\partial_t-L_w$. The linearization $L_w$ has smooth coefficients and belongs to
$\Lcl(V;T,\lambda_1,K,\Lambda_K)$ for every $K$. By Theorem~\ref{thm:linear}, $DG(0)$ is an
isomorphism $\Pp^m_0\to\Pp^{m-1}$.

\begin{proposition}\label{prop:existence}
\lean{ShortTime.EquationData.exists_parabolicSpace_solution, ShortTime.exists_injOn_and_image_mem_nhds, ShortTime.ParabolicSpace.timeDerivExt, ShortTime.ParabolicSpace.existenceMap, ShortTime.LinearOp.isomorphism_eq_heatExt, ShortTime.EquationData.exists_legendre_interval}
\uses{def:data, def:parabolic, def:legendre}
There are $\varepsilon>0$ and $v\in\Pp^m_0$ with the following properties:
\begin{itemize}[nosep]
\item $u=w+v$ has continuous jets lying in $\Par$ for $t\le\varepsilon$;
\item $\partial_tu=F[u]$ on $M\times[0,\varepsilon]$;
\item $u(\cdot,0)=u_0$.
\end{itemize}
Moreover, $v=\Psi(0,f_\varepsilon)$, where $\Psi$ is the map of Proposition~\ref{prop:regularity}
below, taken at the parameter value $0$.
\end{proposition}

\begin{proof}
\uses{lem:formal, lem:cutoff, cor:N, thm:linear, fact:ift, lem:embedding, lem:parabolic-basic, lem:legendre-basic}
\emph{Inverting $G$.} By the inverse function theorem (Fact~\ref{fact:ift}), $G$ restricts to a
$C^\infty$ diffeomorphism from an open $\mathcal U'\ni0$ onto an open $\mathcal V'\ni f$. Shrink
$\mathcal U'$ so that $\sup\abs{j^{2p}v}<\delta_1$ for $v\in\mathcal U'$; this is possible by
Lemma~\ref{lem:embedding}, since $m-1\ge N_0$. By Lemma~\ref{lem:cutoff} with $N=m$ and $m'=m-1$,
$f_\varepsilon\in\mathcal V'$ for $\varepsilon$ small. Put $v=(G|_{\mathcal U'})^{-1}(f_\varepsilon)$.

\emph{The equation on $[0,\varepsilon]$.} Since $f_\varepsilon=0$ for $t\le\varepsilon$, we get
$\partial_tu=F[u]$ almost everywhere on $M\times[0,\varepsilon]$. Both sides are continuous by
Lemma~\ref{lem:embedding}, so equality holds everywhere.

\emph{Initial value and parabolicity.} $u(0)=w(0)=u_0$ by Lemma~\ref{lem:parabolic-basic}(v). The jets
of $u$ lie within $\delta_1$ of those of $w$, hence in $\Par$.
\end{proof}

\subsection{Regularity by Angenent's trick}\label{sec:regularity}

The solution of Proposition~\ref{prop:existence} is only known to lie in $\Pp^m$. We now show that it
is smooth. The idea is to transport it by flows of diffeomorphisms: the transported section solves a
transported equation that depends smoothly on the flow parameter, and the implicit function theorem at
the smooth point $w$ turns this into smooth dependence on the parameter.

\paragraph{Horizontal flows} Let $X$ be a smooth vector field on $M$ with flow $\phi_s$. Parallel
transport along the curves $s\mapsto\phi_s(x)$ defines a one-parameter group $\tilde\phi_s$ of bundle
isometries of $V$ covering $\phi_s$. The pull-back of a section is
\[
  (\tilde\phi_s^*u)(x)=\tilde\phi_s^{-1}\bigl(u(\phi_s(x))\bigr),
\]
and it acts on time-dependent sections slice by slice. Given vector fields $X_1,\dots,X_d$, set
$\tilde\phi_s=\tilde\phi^{X_1}_{s_1}\circ\cdots\circ\tilde\phi^{X_d}_{s_d}$ for $s\in\R^d$.

\begin{lemma}\label{lem:flows}
\lean{ShortTime.HFlow.hpull, ShortTime.HFlow.contMDiffOn_hpullFib, ShortTime.HFlow.jointlySmoothOn_hpull, ShortTime.HFlow.hasDerivAt_hpull_zero, ShortTime.HFlow.hpull_hpull, ShortTime.HFlow.exists_jet_bound_hpull, ShortTime.HFlow.exists_jetTransport, ShortTime.HFlow.natFibreMap, ShortTime.HFlow.natFibreMap_F_zero, ShortTime.HFlow.hpull_nonlinFamily, ShortTime.HFlow.natParam, ShortTime.HFlow.hpull_fibreApply, ShortTime.HFlow.iterPull, ShortTime.HFlow.iterFM, ShortTime.HFlow.iterPull_nonlinFamily, ShortTime.HFlow.iterFlow_chart_patch, ShortTime.HFlow.exists_parabolicInner_hpull_le, ShortTime.HFlow.hpullExt, ShortTime.HFlow.hpullExt_mem_zeroAtStart, ShortTime.HFlow.contDiff_hpullExt_ofSmooth, ShortTime.HFlow.continuous_hpullExt_apply, ShortTime.HFlow.hpullExt_comp_timeDerivExt}
\uses{def:data, def:parabolic, fact:weak}
\begin{enumerate}[label=(\alph*),nosep]
\item For smooth $u$, the map $s\mapsto\tilde\phi_s^*u$ is $C^\infty$ into $\Pp^{m'}$ for every $m'$,
  and
  \[
    \partial_{s_1}\cdots\partial_{s_d}\big|_{s=0}\tilde\phi_s^*u=\nabla_{X_d}\cdots\nabla_{X_1}u .
  \]
\item For $\abs s\le1$ the operators $\tilde\phi_s^*$ are uniformly bounded on $\Pp^{m'}$ and preserve
  $\Pp^{m'}_0$. For every $u\in\Pp^{m'}$ the map $s\mapsto\tilde\phi_s^*u$ is continuous into
  $\Pp^{m'}$.
\item Suppose $u\in\Pp^{m'}$ and $s\mapsto\tilde\phi_s^*u$ is $C^d$ into $\Pp^{m'}$ near $0$. Then the
  weak derivative $\nabla_{X_d}\cdots\nabla_{X_1}u$ exists and equals
  $\partial_{s_1}\cdots\partial_{s_d}|_{0}\tilde\phi_s^*u$, and hence lies in $\Pp^{m'}$.
\item \emph{Naturality.} There are an open set $S\ni0$ and smooth maps $F_s\colon\Dom_s\to V$ for
  $s\in S$, jointly smooth in $(s,t,\zeta)$ on an open set, with $F_0=F$ and $\Dom_0=\Dom$, such that
  \[
    \tilde\phi_s^*\bigl(F[u]\bigr)=F_s\bigl[\tilde\phi_s^*u\bigr]
  \]
  for every smooth $u$ with jets in $\Dom$. Here $F_s[\cdot]$ is formed with the jets of
  $\tilde\phi_s^*u$. Both sides are continuous in $u$ from $\Pp^m_\Dom$ to $\Pp^{m-1}$, and smooth
  sections are dense in $\Pp^m_\Dom$, so the identity extends to all $u\in\Pp^m_\Dom$.
\end{enumerate}
\end{lemma}

\begin{proof}
(a) The derivative of the pull-back by parallel transport along the flow lines of $X$ is $\nabla_X$.
Joint smoothness is standard ODE theory.

(b) Use the chain rule and a change of variables, and approximate by smooth sections.

(c) Pair with a smooth test section $\psi$:
$\ip{\tilde\phi_s^*u}{\psi}=\ip{u}{(\tilde\phi_s^*)^\dagger\psi}$, where
$(\tilde\phi_s^*)^\dagger\psi$ is smooth in $s$. Differentiate at $s=0$ using (a).

(d) Pulling back by $\tilde\phi_s$ intertwines $\nabla$ with the pulled-back connections, and these
differ from $\nabla$ by tensors that depend smoothly on $s$ and vanish at $s=0$. Hence, in a chart,
the jets of $\tilde\phi_s^*u$ are obtained from those of $u$ by a recursion with smooth coefficients.
Covariant jets of sections are not arbitrary (the antisymmetric part of $\nabla^2u$ is a curvature
term), so instead of inverting this jet map we use it in the other direction: there is a smooth family
of fibrewise linear maps $\Upsilon_s(x)\colon\Jet^{2p}_xV\to\Jet^{2p}_{\phi_s(x)}V$, $\Upsilon_0=\id$,
with $\Upsilon_s(x)\bigl(j^{2p}v(x)\bigr)=j^{2p}(\tilde\phi_{-s}^*v)(\phi_s x)$ for all smooth $v$ and
$|s|$ small (glue the chart recursions with a partition of unity). Set
$F_s(t,x,\zeta)=\tilde\phi_s(x)^{-1}F\bigl(t,\phi_s x,\Upsilon_s(x)\zeta\bigr)$, and let $\Dom_s$ be the
corresponding preimage of $\Dom$. Naturality follows from $\tilde\phi_{-s}^*\tilde\phi_s^*=\id$.
\end{proof}

\begin{proposition}[regularity]\label{prop:regularity}
\lean{ShortTime.EquationData.parabolicSpace_solution_smooth, ShortTime.SmoothRep.exists_smooth_rep, ShortTime.SmoothRep.exists_spatialSmooth, ShortTime.SpatialReg.contDiffOn_of_evolution, ShortTime.SmoothRep.evalJet_iterPullExt, ShortTime.Angenent.exists_contDiffAt_of_natural, ShortTime.HFlow.contDiff_iterPullExt_ofSmooth, ShortTime.EquationData.regularity_of_time, ShortTime.EquationData.smoothRep_time, ShortTime.EquationData.exists_iterPull_contDiffOn, ShortTime.Angenent.regularity_core}
\uses{prop:existence, def:solution}
The solution $u$ of Proposition~\ref{prop:existence} is smooth on $M\times[0,\varepsilon]$.
\end{proposition}

\begin{proof}
\uses{lem:flows, fact:ift, fact:weak, lem:spatial-nemytskii, cor:N, fact:sobolev, lem:parabolic-basic}
\emph{Step 1: setting up the implicit function theorem.} Fix smooth vector fields $Y_1,\dots,Y_q$
spanning $TM$ at every point, a length $d\ge1$ and a word $X_1,\dots,X_d$ in the $Y$'s. Near
$(0,0)\in\R^d\times\Pp^m_0$ consider
\[
  H(s,v)=\partial_t\bigl(\tilde\phi_s^*w+v\bigr)-F_s\bigl[\tilde\phi_s^*w+v\bigr]\in\Pp^{m-1}.
\]
By Lemma~\ref{lem:flows}(a),(d) and Corollary~\ref{cor:N}, $H$ is $C^\infty$. Moreover $H(0,\cdot)=G$
and $D_vH(0,0)=DG(0)$ is an isomorphism. By the implicit function theorem (Fact~\ref{fact:ift}) there
are neighbourhoods $S'\times\mathcal V'$ of $(0,f)$ and $\mathcal U''$ of $0$, and a $C^\infty$ map
$\Psi\colon S'\times\mathcal V'\to\mathcal U''$, such that for $(s,y)\in S'\times\mathcal V'$ and
$v\in\mathcal U''$
\[
  H(s,v)=y\iff v=\Psi(s,y).
\]
We choose $\varepsilon$ in Proposition~\ref{prop:existence} so small that
$f_\varepsilon\in\mathcal V'$, and take $v=\Psi(0,f_\varepsilon)$ there. This is consistent with the
proof of that proposition, after shrinking $\mathcal U'$ to lie in $\mathcal U''$.

\emph{Step 2: smooth dependence on the flow parameter.} By Lemma~\ref{lem:flows}(d),
$H(s,\tilde\phi_s^*v)=\tilde\phi_s^*f_\varepsilon$. By Lemma~\ref{lem:flows}(a),(b), for $s$ near $0$
we have $\tilde\phi_s^*v\in\mathcal U''$ and $\tilde\phi_s^*f_\varepsilon\in\mathcal V'$. Therefore
\[
  \tilde\phi_s^*v=\Psi\bigl(s,\tilde\phi_s^*f_\varepsilon\bigr).
\]
The right-hand side is $C^\infty$ in $s$, because $f_\varepsilon$ is smooth. By
Lemma~\ref{lem:flows}(c), $\nabla_{X_d}\cdots\nabla_{X_1}v\in\Pp^m$. Since the word was arbitrary and
$w$ is smooth, Fact~\ref{fact:weak} gives $\partial_t^ju\in L^2(0,\varepsilon;H^k)$ for all $k$ and all
$j\le m$.

\emph{Step 3: regularity in time.} By Lemma~\ref{lem:parabolic-basic}(v), applied with $H^k$ in place
of $L^2$, we get $u\in C([0,\varepsilon];H^k)$ for all $k$. The equation $\partial_tu=F[u]$ says that
\[
  u(t)=u_0+\int_0^t\mathfrak F\bigl(\tau,u(\tau)\bigr)d\tau\qquad\text{in }H^k,
\]
where $\mathfrak F(t,a)=F(t,j^{2p}a)$ is a $C^\infty$ map $H^{k+2p}\to H^k$, jointly in $t$, by
Remark~\ref{lem:spatial-nemytskii}. Hence $u\in C^1([0,\varepsilon];H^k)$ for all $k$. Inductively,
$u\in C^j([0,\varepsilon];H^k)$ for all $j$ and $k$, and $u$ is smooth by Fact~\ref{fact:sobolev}.
\end{proof}

\begin{proof}[Proof of Theorem~\ref{thm:A}]
\lean{ShortTime.EquationData.exists_solution_cov, ShortTime.EquationData.exists_solution_cov'}
\uses{prop:existence, prop:regularity, lem:lowering}
By Lemma~\ref{lem:lowering} we may assume that $V=T^0_sM$. Combine
Propositions~\ref{prop:existence} and~\ref{prop:regularity}, with $T=\varepsilon$.
\end{proof}

\subsection{Uniqueness}\label{sec:uniqueness}

\begin{lemma}[lowering indices]\label{lem:lowering}
\lean{ShortTime.solution_unique_of_covariant, ShortTime.flat, ShortTime.fiberInner_flat, ShortTime.slotPerm, ShortTime.fiberInner_slotPerm, ShortTime.flatCc, ShortTime.unitModel_covGrad_flatCc, ShortTime.covJet_flatCc, ShortTime.contMDiffOn_flat_param, ShortTime.contMDiffOn_flatCoord, ShortTime.trivAt_slotPerm, ShortTime.jointlySmoothOn_flatCc, ShortTime.EquationData.flatData, ShortTime.evalFirst_jetFlat, ShortTime.castIdx_prepend_jetFlat, ShortTime.EquationData.flatData_partialDeriv, ShortTime.EquationData.legendreForm_flatData, ShortTime.EquationData.isParabolicAt_flatData, ShortTime.EquationData.isParabolicSolution_flatData, ShortTime.flatCc_injective, ShortTime.exists_solution_of_covariant, ShortTime.contMDiff_flatSymm, ShortTime.sharpCc, ShortTime.flatCc_sharpCc, ShortTime.EquationData.isParabolicSolution_of_flatData, ShortTime.EquationData.isParabolicSolution_flatData_iff}
\uses{def:data, def:solution, def:legendre}
Let $V=T^r_sM$. Lowering all indices with $g$ is a parallel fibrewise isometry
$\flat\colon T^r_sM\to T^0_{r+s}M$. It carries equation data on $V$ to equation data on $T^0_{r+s}M$,
parabolic solutions to parabolic solutions, and the Legendre condition to the Legendre condition with
the same constant, and conversely through raising indices $\flat^{-1}$. In particular,
Theorems~\ref{thm:A} and~\ref{thm:B} hold on every tensor bundle once they hold on every covariant
tensor bundle $T^0_{s'}M$.
\end{lemma}

\begin{proof}
Since $\flat$ is parallel, $\nabla^k(\flat u)=\sigma_k\,\flat(\nabla^ku)$, where the slot permutation
$\sigma_k$ moves the $k$ slots created by $\nabla^k$ in front of the lowered ones. Hence $\flat$ maps
covariant jets to covariant jets by a fibrewise linear isometry $\flat_J$. Put
$F^\flat(t,\zeta)=\flat F(t,\flat_J^{-1}\zeta)$ on $\Dom^\flat=\flat_J(\Dom)$. Then $u$ solves
$\partial_tu=F[u]$ if and only if $\flat u$ solves $\partial_t(\flat u)=F^\flat[\flat u]$. Finally
$D_{2p}F^\flat=\flat\circ D_{2p}F\circ(\sigma_{2p}\flat)^{-1}$. Prepending covectors commutes with
$\sigma_k\flat$, and evaluating the first $p$ slots of $\sigma_p\flat\xi$ gives $\flat$ of the same
evaluation of $\xi$. Since $\flat$ is an isometry, the Legendre forms of $D_{2p}F$ at $\xi$ and of
$D_{2p}F^\flat$ at $\sigma_p\flat\xi$ agree.
\end{proof}

\begin{lemma}[the difference of two solutions]\label{lem:difference}
\lean{ShortTime.EquationData.exists_linearized, ShortTime.EquationData.F_sub_eq_sum_avg, ShortTime.legendreAt_avgFib, ShortTime.EquationData.exists_avg_coeff, ShortTime.EquationData.contMDiffOn_opFieldOf_avg, ShortTime.EquationData.coordP_avg, ShortTime.EquationData.contMDiffOn_avgIntegrand, ShortTime.EquationData.trivAt_partialDeriv, ShortTime.trivAt_opFieldOf, ShortTime.EquationData.contDiffAt_Fmod, ShortTime.contDiffOn_integral_Icc, ShortTime.contMDiffOn_integral_Icc}
\uses{def:data, def:solution, def:linearization, def:legendre, def:class}
Let $V=T^0_sM$, let $u$ and $v$ be solutions on $[0,T]$, and let $[t_0,t_0+\eta]\subseteq[0,T]$.
Suppose that for $t\in[t_0,t_0+\eta]$ and $\sigma\in[0,1]$ the jets $j^{2p}(v+\sigma(u-v))$ lie in
$\Dom$ and their principal coefficients satisfy the Legendre condition with constant $\lambda$. Then
$h=u-v$ solves $\partial_th=\tilde Lh$ on $[t_0,t_0+\eta]$, where
\[
  \tilde L=\sum_k\tilde A_k\nabla^k,\qquad
  \tilde A_k(t)=\int_0^1D_kF\bigl(t,j^{2p}(v+\sigma h)\bigr)d\sigma,
\]
and $\tilde L$ is a parabolic operator in the sense of Definition~\ref{def:class} with constant
$\lambda$.
\end{lemma}

\begin{proof}
By the fundamental theorem of calculus along the segment of jets,
\[
  F(t,j^{2p}u)-F(t,j^{2p}v)=\sum_k\tilde A_k(\nabla^kh).
\]
The Legendre form is linear and continuous
in the coefficient, so the Legendre condition passes to the average $\tilde A_{2p}$.

The coefficients $\tilde A_k$ are smooth on $M\times[t_0,t_0+\eta]$, up to the boundary. To see this,
fix $x_0\in M$ and work in the trivializations at $x_0$, over an open set $W\ni x_0$. There the
components of $\tilde A_k(x,t)$ are $\int_0^1f(x,t,\sigma)d\sigma$, where $f(x,t,\sigma)$ is the
derivative, in the $k$-th jet slot, of the local representative of $F$ at the coordinate jet of
$v+\sigma h$. This $f$ is smooth on $(W\times[t_0,t_0+\eta])\times[0,1]$. Moreover,
$q\mapsto\int_0^1f(q,\sigma)d\sigma$ is smooth on $W\times S$, for $S$ convex, whenever $f$ is smooth
on $(W\times S)\times[0,1]$: one-sided derivatives are differentiated under the integral sign.
\end{proof}

\begin{lemma}[the Legendre form in coordinates]\label{lem:legendre-coord}
\lean{ShortTime.exists_legendre_coord, ShortTime.legendreForm_expand, ShortTime.continuous_legendreForm_sections, ShortTime.exists_opField_frame, ShortTime.exists_tensor_frame, ShortTime.EquationData.principal_eq_fibMap}
\uses{def:legendre, lem:ibp}
Let $V=T^0_sM$ and $x_0\in M$. Read operators $P\colon(T^0_{s+2p})_x\to(T^0_s)_x$ and tensors
$\xi\in(T^0_{s+p})_x$, for $x$ near $x_0$, in the trivializations at $x_0$. Then, near $x_0$, the
Legendre form $\ip{\xi}{\hatP\xi}$ and the norm $\abs\xi^2$ are continuous functions of $x$ and of
the coordinates of $P$ and $\xi$.
\end{lemma}

\begin{proof}
The Legendre form is linear in $P$ and bilinear in the tensor arguments. Using a bump function,
choose smooth global sections $P_1,\dots,P_N$ and $\Psi_1,\dots,\Psi_K$ whose coordinates near $x_0$
are constant and form bases. In coordinates, $\ip{\xi}{\hatP\xi}$ is then a polynomial in the
coordinates of $P$ and $\xi$, with coefficients $\ip{\Psi_j}{\hatP_i\Psi_l}$. By the pointwise identity
$\ip{\rho\xi}{P^{(p)}\eta}=\ip{\xi}{\hatP\eta}$ from the proof of Lemma~\ref{lem:ibp}, each coefficient
is the pointwise inner product of two smooth sections, hence continuous. The same argument applies to
$\abs\xi^2$.
\end{proof}

\begin{lemma}[local uniqueness]\label{lem:local-unique}
\lean{ShortTime.EquationData.local_unique, ShortTime.EquationData.exists_parabolic_segment, ShortTime.EquationData.tube_local, ShortTime.LegendreAt.mono}
\uses{def:solution, def:legendre, def:class, lem:difference}
Let $V=T^0_sM$, and let $u$ and $v$ be parabolic solutions on $[0,T]$ with $u(t^*)=v(t^*)$ for some
$t^*<T$. Then there is $\eta>0$ with $u=v$ on $[t^*,t^*+\eta]\cap[0,T]$.
\end{lemma}

\begin{proof}
\uses{lem:garding, lem:legendre-coord, fact:gronwall}
\emph{The segment between the jets stays parabolic.} We show that there are $\eta>0$ and $\lambda>0$
such that, for $t\in[t^*,t^*+\eta]$, $x\in M$ and $\sigma\in[0,1]$, the jet
$\zeta=j^{2p}(v+\sigma(u-v))(x,t)$ lies in $\Dom$ and $D_{2p}F(t,\zeta)$ satisfies the Legendre
condition with constant $\lambda$. At $t=t^*$ the segment is the single jet $j^{2p}u(x,t^*)$, because
$u(t^*)=v(t^*)$.

Fix $x_0\in M$ and work in the trivializations at $x_0$. The jet of the segment is continuous in
$(x,t,\sigma)$, and at $t=t^*$ it lies in $\Dom$. Since $\Dom$ is relatively open and $[0,1]$ is
compact, the segment stays in $\Dom$ for $x$ near $x_0$ and $t$ near $t^*$. There the coefficient
$D_{2p}F$ along the segment is continuous in $(x,t,\sigma)$. Let $\lambda_0$ be a Legendre constant of
$u$ at $(t^*,x_0)$, and set
\[
  G(x,t,\sigma,\xi)=(-1)^{p+1}\ip{\xi}{\hatP\xi}-\tfrac{\lambda_0}{2}\abs\xi^2,\qquad
  P=D_{2p}F\bigl(t,j^{2p}(v+\sigma(u-v))(x,t)\bigr),
\]
with $\xi$ in coordinates. By Lemma~\ref{lem:legendre-coord}, $G$ is continuous. At $x=x_0$ and $t=t^*$
it is positive for $\sigma\in[0,1]$ and $\abs\xi=1$. Since $[0,1]$ times the unit sphere is compact,
$G>0$ there for $x$ near $x_0$ and $t\in[t^*,t^*+\eta_{x_0}]$. Both terms of $G$ are quadratic in
$\xi$, so the Legendre condition holds there with constant $\lambda_0/2$. Finitely many of these
neighbourhoods cover $M$; take the least $\eta_{x_0}$ and the least constant.

\emph{The difference solves a linear equation.} Let $h=u-v$. Then
$\partial_th=\tilde Lh=\sum_k\tilde A_k(\nabla^kh)$, where
\[
  \tilde A_k=\int_0^1D_kF\bigl(t,j^{2p}(v+\sigma h)\bigr)d\sigma
\]
is smooth, and $\tilde L$ is parabolic with constant $\lambda$ on $[t^*,t^*+\eta]$, by
Lemma~\ref{lem:difference}.

\emph{Energy estimate.} By Lemma~\ref{lem:garding},
$\frac d{dt}\norm h^2=2\ip h{\tilde Lh}\le2C\norm h^2$. Since $h(t^*)=0$, Gronwall's inequality gives
$h=0$ on $[t^*,t^*+\eta]$.
\end{proof}

\begin{proof}[Proof of Theorem~\ref{thm:B}]
\uses{lem:lowering, lem:local-unique}
By Lemma~\ref{lem:lowering} we may assume that $V=T^0_sM$. Let $u$ and $v$ be parabolic solutions on
$[0,T]$ with the same initial datum. The set of $t\in[0,T]$ with $u(t)=v(t)$ is closed, because $u$
and $v$ are continuous in $t$, and it contains $0$. By Lemma~\ref{lem:local-unique}, each of its
points $t^*<T$ is followed by an interval $[t^*,t^*+\eta]$ inside the set. Since $[0,T]$ is
connected, the set is all of $[0,T]$.
\end{proof}

\subsection{Smooth dependence}

\begin{proof}[Proof of Theorem~\ref{thm:C}]
\uses{thm:linear, cor:N, fact:ift, lem:flows, prop:regularity, lem:parabolic-basic}
\emph{Step 1: the implicit function theorem around $u$.} Let $u$ be a parabolic solution on $[0,T]$,
and set $\bar u=u-Eu_0$. It is smooth with $\bar u(0)=0$, so $\bar u\in\Pp^m_0$. On a neighbourhood of
$(u_0,0)$ in $H^{2pm}(M;V)\times\Pp^m_0$ define
\[
  H(a,v)=\partial_t(Ea+\bar u+v)-F[Ea+\bar u+v].
\]
This is well defined and $C^\infty$ by Lemma~\ref{lem:parabolic-basic}(iv) and Corollary~\ref{cor:N}.
We have $H(u_0,0)=0$, and $D_vH(u_0,0)=\partial_t-L_u$ is an isomorphism by
Theorem~\ref{thm:linear}, applied with the constant
$\kappa_K$ of the compact set of jets of $u$. The implicit function theorem gives $\delta>0$ and a
$C^\infty$ map $a\mapsto v(a)$ with $H(a,v(a))=0$. Put
\[
  \mathcal S(a)=Ea+\bar u+v(a).
\]
Then $\mathcal S(a)$ solves the equation on $[0,T]$ with $\mathcal S(a)(0)=a$. Its jets are uniformly
close to those of $u$, hence in $\Par$, when $\delta$ is small.

\emph{Step 2: smoothness for smooth data.} Let $a$ be smooth. Repeat the proof of
Proposition~\ref{prop:regularity} with the parameters $(a,s)$ and
\[
  H(a,s,v)=\partial_t\bigl(Ea+\tilde\phi_s^*\bar u+v\bigr)-F_s\bigl[Ea+\tilde\phi_s^*\bar u+v\bigr].
\]
Using $\tilde\phi_s^*E=E\tilde\phi_s^*$ one gets
$\tilde\phi_s^*v(a)=\Psi(\tilde\phi_s^*a,s)$ for small $s$. This is $C^\infty$ in $s$, because
$s\mapsto\tilde\phi_s^*a$ is $C^\infty$ into $H^{2pm}$. Hence $\mathcal S(a)$ is smooth, arguing as in
Steps~2--3 of that proof on all of $[0,T]$.
\end{proof}

\begin{proof}[Proof of Corollary~\ref{cor:C}]
\uses{thm:C, thm:B, lem:embedding}
\emph{Remark on the formalization.} The Lean proof does not use the smooth map $\mathcal S$ of
Theorem~\ref{thm:C}. It perturbs the fixed smooth solution $u$ by the time-independent family
$a-u_0$, applies the implicit function theorem with parameter, and uses regularity in trace form at
the fixed base solution (see \texttt{blueprint/DECISIONS.md}). The argument below is the original route.
\begin{itemize}[nosep]
\item Consistency across levels: for $m'\ge m$ and smooth $a$ in both domains, the level-$m'$ and
  level-$m$ solutions coincide by Theorem~\ref{thm:B}.
\item If $u_0^{(k)}\to u_0$ in $C^\infty$, then for every $m$ we have
  $\norm{u_0^{(k)}-u_0}_{H^{2pm}}<\delta_m$ for large $k$.
\item Therefore $u^{(k)}=\mathcal S_m(u_0^{(k)})\to u$ in $\Pp^m$ for every $m$, and hence in
  $C^\infty$ by Lemma~\ref{lem:embedding}.
\end{itemize}
\end{proof}

% ===========================================================================
\section{Examples}\label{sec:examples}
% ===========================================================================

In each example we identify the principal coefficient and check the Legendre condition using
Lemma~\ref{lem:legendre-examples}. The geometric conclusions that go beyond short-time existence are
not discussed; an example is the invariance of a target manifold, or the passage from a gauge-fixed
flow to the original flow.

\subsection{Quasilinear equations of product type}
This is the setting of Corollary~\ref{cor:MM}. It includes the $Q$-curvature flow on four-manifolds,
whose principal part is $-e^{-4w}\Delta_{g_0}^2$ (see \cite{Brendle,MalchiodiStruwe}), and the
Willmore-type flows of \cite{KuwertSchaetzle,Polden}, once these are written as scalar equations for a
graph function. \todo{Check the exact forms against the references.}

\subsection{Polyharmonic heat flows of maps}
Let $N\subset\R^K$ be a closed submanifold, and let $V=M\times\R^K$ be the trivial bundle with the flat
connection. The extrinsic polyharmonic map heat flow of order $2p$ has the form
\[
  \partial_tu=-(\Lap)^pu+\mathcal Q\bigl(u,\nabla u,\dots,\nabla^{2p-1}u\bigr),
\]
with $\mathcal Q$ smooth for $u$ in a tubular neighbourhood of $N$; see \cite{EellsSampson} for $p=1$
and \cite{Lamm,Gastel} for $p\ge2$. \todo{Verify that the lower-order structure is as stated for all
$p$.} By Lemma~\ref{lem:legendre-examples}(b), $\Par=\Dom$ is the set of jets whose zeroth component
lies in the tubular neighbourhood. Theorems~\ref{thm:A}--\ref{thm:C} give short-time existence,
uniqueness and smooth dependence for every $p$, including the sixth-order ($p=3$) triharmonic flow.

\subsection{The Ricci--DeTurck flow}\label{ex:rdt}
Fix a background metric $\tilde g$ on $M$, and use it as the metric $g$ of Section~\ref{sec:setting}.
Let $V=\operatorname{Sym}^2T^*M$ with the induced metric and connection. The Ricci--DeTurck flow
\cite{DeTurck,Hamilton82} can be written as
\[
  \partial_tg_{ij}=(g^{-1})^{ab}\tilde\nabla_a\tilde\nabla_bg_{ij}+\mathcal Q(g,g^{-1},\tilde\nabla g,\widetilde{\mathrm{Rm}})_{ij},
\]
with $\mathcal Q$ smooth; see \cite{Shi,ChowKnopf}. Here $\Dom=\Par$ is the set of jets whose zeroth
component is positive definite. The principal coefficient is $P(X)=(\zeta_0^{-1})^{ab}X_{ab}$, and
$\kappa(P)$ is the smallest eigenvalue of $\zeta_0^{-1}$ relative to $\tilde g$, by
Lemma~\ref{lem:legendre-examples}(a) with $p=1$. Theorems~\ref{thm:A}--\ref{thm:C} therefore give
short-time existence, uniqueness and smooth dependence for the Ricci--DeTurck flow. The Ricci flow is
recovered by DeTurck's diffeomorphisms, which is a separate ODE argument.

\subsection{Parabolic Monge--Amp\`ere equations}
A fully nonlinear example with a nontrivial parabolicity domain is
\[
  \partial_tu=\log\frac{\det\bigl(g+\nabla^2u\bigr)}{\det g}+h(x,t,u,\nabla u),
\]
for scalar $u$. Take $\Dom=\Par$ to be the set of jets with $g+\operatorname{Sym}\zeta_2>0$. The
principal coefficient is $P(X)=\bigl((g+\operatorname{Sym}\zeta_2)^{-1}\bigr)^{ab}X_{ab}$, which
satisfies the Legendre condition on $\Dom$.

The K\"ahler--Ricci flow written as a parabolic complex Monge--Amp\`ere equation for a potential
\cite{Cao} has the same structure. On a K\"ahler manifold, $\sqrt{-1}\partial\bar\partial\varphi$ is a
linear function of $\nabla^2\varphi$, and the principal coefficient is a positive multiple of
$g_\varphi^{-1}$ acting on the $J$-invariant part. \todo{Write out the K\"ahler case in detail.}

% ===========================================================================
\section{Remarks on the scope}\label{sec:scope}
% ===========================================================================

\begin{enumerate}[label=(\arabic*)]
\item \emph{Gauge-degenerate flows.} For the Ricci flow, the mean curvature and Willmore flows of
  immersions, the Yang--Mills flow and similar equations, $F$ is degenerate in the directions of a gauge
  group, so the Legendre condition fails. The standard remedy (DeTurck's trick) is to add a
  gauge-fixing term, solve the resulting strictly parabolic equation, which is covered by our results,
  and then undo the gauge by solving an ODE. A general version of this last step, stated abstractly
  for flows equivariant under a gauge group, would be a natural companion to the present results.
\item \emph{Strong parabolicity without a Legendre representation.} By Remark~\ref{rem:sos}, such
  operators exist in general. They require G\aa rding's inequality in its general form (localization
  and Fourier analysis). The rest of the argument is unchanged, because everything downstream of
  Lemma~\ref{lem:garding} uses only G\aa rding's inequality, and uses it only for operators
  $L^{(r)}$ obtained by differentiating (Lemma~\ref{lem:commute}).
\item \emph{Nonlocal lower-order terms.} Suppose $F$ depends in addition on finitely many
  integrals of the form $\int_M\Phi(t,j^{2p-1}u)\dmu$, such as normalizations or averages. Then
  Proposition~\ref{prop:nemytskii} and the rest of the argument go through with trivial changes.
\item \emph{Analytic dependence.} If $F$ is real analytic in $\zeta$, the maps in the implicit function
  theorem are analytic, and $\mathcal S$ in Theorem~\ref{thm:C} is analytic.
\item \emph{General bundles versus tensor bundles.} The paper is written for arbitrary Riemannian
  bundles. The formal statements treat the tensor bundles $V=T^r_sM$ of type $(r,s)$, with the
  metric and connection induced by $g$ and its Levi-Civita connection. These are closed under the
  operation $W\mapsto T^*M\otimes W$ used in Lemma~\ref{lem:commute}, and for them the horizontal
  flows of Section~\ref{sec:regularity} can be replaced by natural pull-backs of tensors.
  Vector-valued systems, $V=T^r_sM\otimes\R^K$ with $K\ge2$, are not covered by the formal
  statements; they need only an extra fixed factor in the model fibre.
\end{enumerate}

% ===========================================================================
\section{Notes on the formalization}\label{sec:formal}
% ===========================================================================

\subsection{Dependency structure}

Every numbered item of this paper carries two annotations for software, invisible in the printed
version: the Lean declarations formalizing it, and the items its statement and proof use.
A script in the repository turns them into a dependency graph with formalization status, which it
writes to \texttt{blueprint/map.md}. It can also \emph{pin} approved Lean statements, so
that later changes to them are detected. At the coarsest level the proof has four layers, each of
which uses only the ones before it.

\begin{enumerate}[label=\textbf{L\arabic*}]
\item \emph{Geometry and algebra:} Lemma~\ref{lem:ricci}, Definition~\ref{def:legendre},
  Lemmas~\ref{lem:legendre-basic} and~\ref{lem:legendre-examples}.
\item \emph{Function spaces:} Lemmas~\ref{lem:sobolev-basic}, \ref{lem:interp},
  \ref{lem:parabolic-basic} and~\ref{lem:embedding}, Proposition~\ref{prop:nemytskii},
  Corollary~\ref{cor:N}, and the imported
  Facts~\ref{fact:sobolev}, \ref{fact:weak}, \ref{fact:eigen} and~\ref{fact:elliptic}.
\item \emph{Linear theory:} Lemmas~\ref{lem:ibp}, \ref{lem:garding}, \ref{lem:commute}
  and~\ref{lem:continuity}, Propositions~\ref{prop:energy}, \ref{prop:apriori}
  and~\ref{prop:model}, and Theorem~\ref{thm:linear}.
\item \emph{Nonlinear theory:} Lemmas~\ref{lem:formal}, \ref{lem:cutoff} and~\ref{lem:flows},
  Propositions~\ref{prop:existence} and~\ref{prop:regularity}, Theorems~\ref{thm:A},
  \ref{thm:B} and~\ref{thm:C}, and the imported Fact~\ref{fact:ift}.
\end{enumerate}

Layers L3 and L4 interact only through Theorem~\ref{thm:linear}, and only for operators with smooth
coefficients. The two layers can be developed in parallel once that statement is fixed.

\subsection{The trusted statements}

Theorems~\ref{thm:A} and~\ref{thm:B}, and Corollary~\ref{cor:C} (the sequential form of
Theorem~\ref{thm:C}), mention only
smooth sections, covariant derivatives, time derivatives and the pointwise algebraic
Definition~\ref{def:legendre}. These are the statements that should appear in the final formal
version, so that the definitions a reader has to trust stay small. The parabolic Sobolev spaces are
internal to the proofs.

\subsection{Imported facts and candidate sources}

The following table records where each imported or reusable ingredient might come from. The entries
marked ``to check'' are based on file names and statements seen in the library, not on a build.

\medskip
\begin{center}
\footnotesize
\begin{tabular}{@{}p{4.1cm}p{2.9cm}p{6.4cm}@{}}
\toprule
Ingredient & Used in & Candidate source\\
\midrule
Ricci identity, commutators & Lemma~\ref{lem:ricci} & \cite{DGlib} \texttt{Tensor/RicciIdentity} (partial)\\
Divergence theorem, integration by parts & throughout & \cite{DGlib} \texttt{Analysis/Integration}\\
Iterated covariant derivatives & \S\ref{sec:setting} & \cite{DGlib} \texttt{iteratedCovGrad}\\
Interpolation (log-convexity) & Lemma~\ref{lem:interp} & \cite{DGlib} \texttt{Sobolev/GagliardoNirenberg/\allowbreak FiberNorm} (to check)\\
Sobolev embedding & Fact~\ref{fact:sobolev} & \cite{DGlib} \texttt{Sobolev/Embedding/Tensor}\\
Weak $=$ strong derivatives & Fact~\ref{fact:weak} & \cite{DGlib} \texttt{Sobolev/Intrinsic} (to check)\\
Eigenbasis of $\Lap$ & Fact~\ref{fact:eigen} & \cite{DGlib} \texttt{Spectral/Tensor} (tensor bundles)\\
Elliptic estimate & Fact~\ref{fact:elliptic} & \cite{DGlib} \texttt{Spectral/Tensor/\allowbreak SobolevScale/Jet} (to check)\\
$H^1(0,T;X)\subset C([0,T];X)$ & Lemma~\ref{lem:parabolic-basic} & \cite{DGlib} \texttt{Parabolic/TimeSobolev/H1}\\
Per-mode ODE solution & Proposition~\ref{prop:model} & \cite{DGlib} \texttt{Parabolic/\allowbreak MaximalRegularity/PerMode}\\
Flows of vector fields & Lemma~\ref{lem:flows} & \cite{DGlib} \texttt{Analysis/ODE}, Mathlib\\
Inverse/implicit function theorem & Fact~\ref{fact:ift} & Mathlib\\
Gronwall's inequality & Prop.~\ref{prop:energy}, Thm.~\ref{thm:B} & Mathlib\\
Method of continuity & Lemma~\ref{lem:continuity} & new (short)\\
\bottomrule
\end{tabular}
\end{center}

\subsection{Points in this draft that need checking}

\begin{itemize}[nosep]
\item Remark~\ref{rem:sos}: the precise statement and the references.
\item Section~\ref{sec:examples}: the exact form of each example equation against the literature.
\item Lemma~\ref{lem:flows} and Step~3 of Proposition~\ref{prop:regularity}: these are written at the
  level of ``routine'', and a formalization will need them spelled out.
\item The values of the constants $K$ required in Propositions~\ref{prop:energy}
  and~\ref{prop:apriori} are generous rather than sharp; any finite value suffices.
\end{itemize}

% ===========================================================================
\appendix
\section{Abstract tools}\label{app:abstract}
% ===========================================================================

\begin{fact}[implicit function theorem {\cite{Lang}}]\label{fact:ift}
\source{Mathlib: HasStrictFDerivAt.implicitFunction, ContDiffAt.to_localInverse}
Let $X$, $Y$ and $Z$ be Banach spaces, $U\subseteq Z\times X$ open, and $H\colon U\to Y$ of class
$C^k$ with $k\ge1$. Suppose $H(z_0,x_0)=y_0$ and $D_xH(z_0,x_0)\colon X\to Y$ is invertible. Then
there are open neighbourhoods $Z'\ni z_0$, $Y'\ni y_0$ and $X'\ni x_0$, and a $C^k$ map
$\Psi\colon Z'\times Y'\to X'$, such that for $(z,y,x)\in Z'\times Y'\times X'$
\[
  H(z,x)=y\iff x=\Psi(z,y).
\]
Taking $Z$ to be a point gives the inverse function theorem.
\end{fact}

\begin{fact}[Gronwall]\label{fact:gronwall}
\source{Mathlib: norm_le_gronwallBound_of_norm_deriv_right_le}
Let $\varphi\in C^1([t_0,t_1])$ with $\varphi(t_0)=0$ and $\varphi'\le a\varphi+b$, where $a\ge0$ and
$b\ge0$ is continuous. Then $\varphi(t)\le e^{a(t-t_0)}\int_{t_0}^tb$.
\end{fact}

% ===========================================================================
\begin{thebibliography}{99}

\bibitem{Amann} H. Amann, \emph{Linear and Quasilinear Parabolic Problems. Vol.~I: Abstract Linear
Theory}, Monographs in Mathematics 89, Birkh\"auser, 1995.

\bibitem{Angenent90} S. B. Angenent, Nonlinear analytic semiflows, \emph{Proc. Roy. Soc. Edinburgh
Sect.~A} \textbf{115} (1990), 91--107.

\bibitem{Aubin} T. Aubin, \emph{Some Nonlinear Problems in Riemannian Geometry}, Springer, 1998.

\bibitem{Baker} C. Baker, \emph{The mean curvature flow of submanifolds of high codimension}, PhD
thesis, Australian National University, 2010.

\bibitem{Brendle} S. Brendle, Global existence and convergence for a higher order flow in conformal
geometry, \emph{Ann. of Math.} \textbf{158} (2003), 323--343.

\bibitem{Cao} H.-D. Cao, Deformation of K\"ahler metrics to K\"ahler--Einstein metrics on compact
K\"ahler manifolds, \emph{Invent. Math.} \textbf{81} (1985), 359--372.

\bibitem{ChowKnopf} B. Chow and D. Knopf, \emph{The Ricci Flow: An Introduction}, Mathematical Surveys
and Monographs 110, AMS, 2004.

\bibitem{DeTurck} D. M. DeTurck, Deforming metrics in the direction of their Ricci tensors, \emph{J.
Differential Geom.} \textbf{18} (1983), 157--162.

\bibitem{DGlib} \texttt{differential-geometry}: a Lean 4 library for differential geometry and
geometric analysis, \url{https://github.com/qinz1yang/differential-geometry}, version 0.1.3, 2026.

\bibitem{EellsSampson} J. Eells and J. H. Sampson, Harmonic mappings of Riemannian manifolds,
\emph{Amer. J. Math.} \textbf{86} (1964), 109--160.

\bibitem{Eidelman} S. D. Eidelman, \emph{Parabolic Systems}, North-Holland, 1969.

\bibitem{Gastel} A. Gastel, The extrinsic polyharmonic map heat flow in the critical dimension,
\emph{Adv. Geom.} \textbf{6} (2006), 501--521. \todo{check}

\bibitem{GT} D. Gilbarg and N. S. Trudinger, \emph{Elliptic Partial Differential Equations of Second
Order}, Springer, 2001 (reprint of the 1998 edition).

\bibitem{Hamilton75} R. S. Hamilton, \emph{Harmonic Maps of Manifolds with Boundary}, Lecture Notes in
Mathematics 471, Springer, 1975.

\bibitem{Hamilton82} R. S. Hamilton, Three-manifolds with positive Ricci curvature, \emph{J.
Differential Geom.} \textbf{17} (1982), 255--306.

\bibitem{Hebey} E. Hebey, \emph{Nonlinear Analysis on Manifolds: Sobolev Spaces and Inequalities},
Courant Lecture Notes 5, AMS, 2000.

\bibitem{Hilbert} D. Hilbert, \"Uber die Darstellung definiter Formen als Summe von Formenquadraten,
\emph{Math. Ann.} \textbf{32} (1888), 342--350.

\bibitem{HuiskenPolden} G. Huisken and A. Polden, Geometric evolution equations for hypersurfaces, in:
\emph{Calculus of Variations and Geometric Evolution Problems (Cetraro, 1996)}, Lecture Notes in
Mathematics 1713, Springer, 1999, 45--84.

\bibitem{KuwertSchaetzle} E. Kuwert and R. Sch\"atzle, Gradient flow for the Willmore functional,
\emph{Comm. Anal. Geom.} \textbf{10} (2002), 307--339.

\bibitem{Lamm} T. Lamm, Heat flow for extrinsic biharmonic maps with small initial energy, \emph{Ann.
Global Anal. Geom.} \textbf{26} (2004), 369--384. \todo{check}

\bibitem{Lang} S. Lang, \emph{Fundamentals of Differential Geometry}, Graduate Texts in Mathematics
191, Springer, 1999.

\bibitem{LionsMagenes} J.-L. Lions and E. Magenes, \emph{Non-Homogeneous Boundary Value Problems and
Applications, Vol.~I}, Springer, 1972.

\bibitem{Lunardi} A. Lunardi, \emph{Analytic Semigroups and Optimal Regularity in Parabolic Problems},
Birkh\"auser, 1995.

\bibitem{MalchiodiStruwe} A. Malchiodi and M. Struwe, $Q$-curvature flow on $S^4$, \emph{J.
Differential Geom.} \textbf{73} (2006), 1--44.

\bibitem{MM} C. Mantegazza and L. Martinazzi, A note on quasilinear parabolic equations on manifolds,
\emph{Ann. Sc. Norm. Super. Pisa Cl. Sci. (5)} \textbf{XI} (2012), 857--874.

\bibitem{mathlib} The mathlib Community, The Lean mathematical library, in: \emph{Proceedings of the
9th ACM SIGPLAN International Conference on Certified Programs and Proofs (CPP 2020)}, 367--381.

\bibitem{Motzkin} T. S. Motzkin, The arithmetic-geometric inequality, in: \emph{Inequalities}
(O. Shisha, ed.), Academic Press, 1967, 205--224.

\bibitem{Polden} A. Polden, \emph{Curves and surfaces of least total curvature and fourth-order
flows}, PhD thesis, Universit\"at T\"ubingen, 1996.

\bibitem{Sharples} J. J. Sharples, Linear and quasilinear parabolic equations in Sobolev space,
\emph{J. Differential Equations} \textbf{202} (2004), 111--142.

\bibitem{Shi} W.-X. Shi, Deforming the metric on complete Riemannian manifolds, \emph{J. Differential
Geom.} \textbf{30} (1989), 223--301.

\bibitem{Solonnikov} V. A. Solonnikov, On boundary value problems for linear parabolic systems of
differential equations of general form, \emph{Trudy Mat. Inst. Steklov.} \textbf{83} (1965), 3--163.

\end{thebibliography}

\end{document}
